\documentclass[11pt,a4paper]{article}

% ── Packages ─────────────────────────────────────────────────────────────────
\usepackage[utf8]{inputenc}
\usepackage[T1]{fontenc}
\usepackage{lmodern}
\usepackage[margin=2.5cm]{geometry}
\usepackage{amsmath,amssymb,amsthm}
\usepackage{booktabs}
\usepackage{listings}
\usepackage{xcolor}
\usepackage{hyperref}
\usepackage{url}
\usepackage{graphicx}
\usepackage{microtype}
\usepackage{parskip}
\usepackage{authblk}
\usepackage{abstract}

% ── Theorem environments ──────────────────────────────────────────────────────
\newtheorem{theorem}{Theorem}[section]
\newtheorem{lemma}[theorem]{Lemma}
\newtheorem{definition}[theorem]{Definition}
\newtheorem{remark}[theorem]{Remark}
\newtheorem{corollary}[theorem]{Corollary}

% ── Code listing style ────────────────────────────────────────────────────────
\lstset{
  basicstyle=\ttfamily\small,
  breaklines=true,
  frame=single,
  backgroundcolor=\color{gray!5},
  keywordstyle=\bfseries,
  numbers=left,
  numberstyle=\tiny\color{gray},
}

% ── Metadata ──────────────────────────────────────────────────────────────────
\title{\textbf{NAPQES: A Single-Primitive Authenticated Encryption Scheme\\
Built from HMAC-SHA256}}

\author[1]{Abdeslam Jakjoud}
\author[1]{Karim Amor}
\affil[1]{QAegis, France}
\date{July 2026 \\ \small{Version 2 --- submitted to IACR ePrint (revises Version 1, May 2026)}}

% ─────────────────────────────────────────────────────────────────────────────
\begin{document}
\maketitle

% ── Abstract ─────────────────────────────────────────────────────────────────
\begin{abstract}
We present NAPQES, an Authenticated Encryption with Associated Data (AEAD)
scheme whose security reduces exclusively to the pseudorandomness of
HMAC-SHA256.  Unlike AES-GCM and ChaCha20-Poly1305, NAPQES does not employ
finite-field arithmetic, group operations, or AES S-box permutations, and is
therefore not subject to Shor's algorithm or hidden-subgroup quantum attacks.
The construction combines a prime-indexed token cipher, HMAC-derived noise
injection to achieve ciphertext pseudorandomness, HMAC-derived plaintext
padding for cross-implementation determinism, and a standard HMAC-SHA256
authentication tag.
\textbf{As of this version, the primary and recommended NAPQES construction
is v8}: it keys the arithmetic (token) layer and the HMAC domain-derivation
layer with two \emph{independently sampled} secrets, $\mathbf{k}$ (a prime
tuple) and $\mathsf{sk}$ (a uniform 256-bit string), and derives the nonce
synthetically as $N=\mathrm{HMAC}(\mathsf{sk}, \texttt{0x0A}\|A\|M)$ in the
manner of AES-GCM-SIV, rather than sampling it independently at random. This
single change closes three findings that apply to every earlier wire format
(v1--v7): exact key recovery under nonce reuse, a non-standard
low-entropy-key PRF hypothesis, and a security-reduction simulation gap.
We prove IND-CPA, INT-CTXT, and IND-CCA security for v8 under the standard,
uniform-key PRF assumption on HMAC-SHA256 (Theorems~\ref{thm:ind-cpa-v8},
\ref{thm:int-ctxt}, \ref{thm:ind-cca}), with no key-entropy or
key-guessing term, and with nonce misuse degrading gracefully to disclosed
message-equality leakage instead of key recovery. The legacy, single-secret
v7 wire format is retained, fully specified, and proved secure under a
weaker, explicit key-entropy assumption
(Theorem~\ref{thm:ind-cpa}, Remark~\ref{rem:min-K}) for backward
compatibility only; it is no longer recommended for new deployments.
We discuss the post-quantum security level, and provide performance
measurements and implementation notes. The streaming Release of Unverified
Plaintext (RUP) caveat present in earlier versions is addressed by an
online-AE streaming format
(\texttt{encrypt\_stream\_ae}/\texttt{decrypt\_stream\_ae}) that verifies a
per-chunk HMAC-SHA256 tag before releasing any plaintext to the caller.
A reference implementation in Rust (v8) and Python/Rust/C (v7), together
with a 30-vector Known-Answer Test corpus and a NIST SP 800-22 Rev 1a
statistical analysis, are published at \url{https://github.com/[REPO]}.
\end{abstract}

\paragraph{Changes since Version 1 (May 2026).} This version:
\begin{enumerate}
  \item promotes the v8 independent-key-schedule, synthetic-nonce
        construction (Section~\ref{sec:v8}) from an appendix-style fix to
        the paper's primary, recommended construction, and states
        standalone IND-CPA/INT-CTXT/IND-CCA theorems for it
        (Theorem~\ref{thm:ind-cpa-v8} and Corollary~\ref{cor:v8-security});
  \item retitles the single-secret construction analysed in
        Sections~\ref{sec:algorithm-triple}--\ref{sec:security} as the
        \emph{legacy v7 variant}, retained for backward compatibility and
        historical completeness, and makes explicit that it is no longer
        the recommended default;
  \item updates the abstract, comparison table
        (Table~\ref{tab:comparison}), and conclusion to reflect v8 as the
        default; no proof, theorem statement, or wire-format byte layout
        from Version 1 is altered --- every Version 1 result is preserved
        verbatim and continues to apply to the legacy v7 variant it was
        proved against.
\end{enumerate}

\tableofcontents
\clearpage

% ─────────────────────────────────────────────────────────────────────────────
\section{Introduction}
\label{sec:intro}
% ─────────────────────────────────────────────────────────────────────────────

Modern AEAD standards—AES-GCM~\cite{nist-gcm} and
ChaCha20-Poly1305~\cite{rfc7539}—have achieved widespread deployment.
Their security, however, ultimately depends on algebraic structures
(finite-field multiplication for GCM, add-rotate-XOR chains for ChaCha20)
whose resilience against quantum adversaries running variants of Shor's
algorithm is an active research question~\cite{nistpq2022}.

NAPQES takes a more conservative approach: the \emph{only} cryptographic
primitive in the construction is HMAC-SHA256, approved under NIST FIPS 198-1
and FIPS 180-4.  The security of the entire scheme therefore reduces to the
pseudorandomness of HMAC-SHA256—a well-studied assumption with no known
quantum speedup beyond Grover's $O(\sqrt{N})$ search, which reduces the
effective security level from 256~bits to approximately 128~bits while
remaining far above all recommended thresholds.

\paragraph{Contributions.}
\begin{enumerate}
  \item We specify \textbf{v8}, the primary and recommended NAPQES
        construction (Section~\ref{sec:v8}): an independent
        HMAC-domain-derivation secret $\mathsf{sk}$ and arithmetic-layer
        key $\mathbf{k}$, sampled independently, together with a synthetic,
        message-bound nonce $N=\mathrm{HMAC}(\mathsf{sk},\texttt{0x0A}\|A\|M)$.
        We prove IND-CPA, INT-CTXT, and IND-CCA security for v8
        (Theorem~\ref{thm:ind-cpa-v8} and Corollary~\ref{cor:v8-security})
        under the \emph{standard, uniform-key} PRF assumption on
        HMAC-SHA256, with no key-entropy term, and show that nonce reuse
        degrades gracefully to disclosed plaintext-equality leakage rather
        than exact key recovery.
  \item We fully specify the shared token/wire-format machinery
        (Section~\ref{sec:construction}) --- key serialisation, HMAC
        derivation schedule, token construction, padding scheme, and
        authentication tag --- common to both v8 and the legacy v7 wire
        format.
  \item For backward compatibility, we retain the legacy, single-secret
        v7 construction and its full security analysis
        (Sections~\ref{sec:algorithm-triple}--\ref{sec:security}): IND-CPA
        security under the PRF assumption on HMAC-SHA256, with an explicit
        key-entropy term accounting for the fact that its HMAC key is a
        low-entropy, structured prime tuple rather than a uniform
        bit-string (Theorem~\ref{thm:ind-cpa}, Section~\ref{sec:cvf8-fix}),
        and show that the authentication tag provides EUF-CMA integrity. v7
        is no longer recommended for new deployments (Section~\ref{sec:v8}).
  \item We provide a post-quantum analysis (Section~\ref{sec:pq}) and a
        comparison with AES-GCM and ChaCha20-Poly1305
        (Section~\ref{sec:comparison}).
  \item We report NIST SP 800-22 Rev 1a statistical test results on 10 million
        bits of NAPQES ciphertext output (Section~\ref{sec:sts}), a
        30-vector KAT corpus (Section~\ref{sec:kats}), and a fuzz harness
        covering the full decoder pipeline (Section~\ref{sec:fuzz}).
  \item We discuss known caveats (Section~\ref{sec:caveats}): the 16-bit
        plaintext length cap (CAV-002), the padding length-bucket leakage
        (CAV-003), and the ciphertext expansion bound (CAV-004).  The
        streaming Release of Unverified Plaintext issue (CAV-001) has been
        resolved by the online-AE streaming format introduced in this version.
  \item We state the normative status of, and selection rule among, the
        three wire encodings — block-mode (v7/v8), the deprecated basic
        streaming format, and the online-AE streaming format
        (Section~\ref{sec:format-applicability}) — closing audit finding
        CVF7.
\end{enumerate}

% ─────────────────────────────────────────────────────────────────────────────
\section{Related Work}
\label{sec:related}
% ─────────────────────────────────────────────────────────────────────────────

\paragraph{AES-GCM~\cite{nist-gcm}.}
AES-GCM combines AES-CTR for confidentiality with a GHASH polynomial MAC
over $\mathrm{GF}(2^{128})$.  The Forbidden Attack~\cite{joux2006} and
nonce-reuse attacks~\cite{joux-nonce} demonstrate that the algebraic structure
of GHASH can be exploited when the polynomial MAC key is recovered (e.g.,
under nonce misuse), algebraically recovering the authentication key.
NAPQES uses HMAC-SHA256 for authentication, which does not admit the
Forbidden Attack.

\paragraph{ChaCha20-Poly1305~\cite{rfc7539}.}
ChaCha20 is an ARX cipher whose security is not reduced to a well-studied
hardness assumption; Poly1305 is again a polynomial MAC, subject to the same
class of Forbidden Attacks as GHASH when the one-time key is reused.  Neither
ChaCha20 nor Poly1305 is approved under current FIPS standards.

\paragraph{Ascon~\cite{ascon2021}.}
Ascon (NIST LWC winner) is a sponge-based AEAD targeting lightweight
environments.  Its security is based on the hardness of distinguishing the
Ascon permutation from a random permutation, which involves xor-and-rotate
operations over 64-bit words.  Ascon is not yet in a FIPS-approved family.

\paragraph{HMAC-based constructions.}
HKDF~\cite{rfc5869} and HMAC-DRBG~\cite{nist-drbg} demonstrate that
HMAC-SHA256 is a sufficient basis for key derivation and pseudorandom
generation.  NAPQES extends this philosophy to the full AEAD context.

% ─────────────────────────────────────────────────────────────────────────────
\section{Construction}
\label{sec:construction}
% ─────────────────────────────────────────────────────────────────────────────

\begin{remark}[Reading guide: v8 is the recommended default]
\label{rem:reading-guide}
Sections~\ref{sec:construction}.1--\ref{sec:construction}.9 below (Notation
through Wire Format Version 7) specify the token/wire-format machinery ---
domain-separated HMAC derivation, noise probability, padding, token
construction, and the masked-blob wire format --- shared verbatim by both
NAPQES key schedules described in this paper. Section~\ref{sec:algorithm-triple}
then instantiates that machinery as the formal algorithm triple
\emph{under the legacy, single-secret v7 key schedule}, in which one secret
$\mathbf{k}$ keys both the arithmetic layer and every HMAC domain
derivation; Section~\ref{sec:v8} instantiates the identical machinery under
an \emph{independent} pair of secrets $(\mathbf{k},\mathsf{sk})$ with a
synthetic nonce. \textbf{Section~\ref{sec:v8} (v8) is the primary,
recommended construction as of this version}; Section~\ref{sec:algorithm-triple}
(v7) is retained, fully specified, and separately proved secure, for
backward compatibility and historical completeness only. Readers
interested only in the recommended construction may read
Sections~\ref{sec:construction}.1--\ref{sec:construction}.9 for the shared
machinery and then skip directly to Section~\ref{sec:v8}.
\end{remark}

\subsection{Notation}

Throughout this paper:
\begin{itemize}
  \item $\|$ denotes byte concatenation.
  \item $\mathrm{HMAC}(K, M)$ denotes HMAC-SHA256 keyed with $K$ over message $M$.
  \item $\mathrm{be5}(x)$ denotes the 5 least-significant bytes of integer $x$,
        big-endian encoded.
  \item $\mathrm{be8}(x)$ denotes the 8-byte big-endian encoding of integer $x$
        (the fixed-width token field introduced by the CVF1 fix; see
        Section~\ref{sec:construction} "Wire Format (Version 7)").
  \item $\mathrm{varint}(x)$ denotes unsigned LEB128 encoding of non-negative
        integer $x$, defined precisely in Definition~\ref{def:leb128} below.
        \textbf{Retired for block-mode tokens as of the CVF1 fix}
        (Section~\ref{sec:construction}); still used by the streaming format
        and by legacy v6 ciphertext decoding.
  \item $\mathcal{P}$ denotes the set of prime integers in $[10^6, 9.9 \times 10^6]$.
\end{itemize}

\begin{remark}[Audit finding CVF17: LEB128 used without a definition or normative reference]
\label{rem:cvf17}
The notation above previously stated only that $\mathrm{varint}(x)$ is
``unsigned LEB128 encoding of non-negative integer $x$,'' without ever
defining LEB128 or citing a normative reference for it. As written, the
byte layout, the continuation-bit convention, and whether a canonical
(minimal-length) encoding is required were all left implicit, so the wire
format could not be reproduced from the specification alone. Beyond being
undefined, variable-length LEB128 decoding is error-prone in ways that
have security relevance for a decoder (overlong/non-canonical encodings,
continuation-bit handling, unbounded length); those concrete decoding
hazards are tracked as a separate implementation-level finding and are not
the subject of this remark. Definition~\ref{def:leb128} below states the
encoding explicitly and specifies canonicality and rejection rules for
decoding.
\end{remark}

\begin{definition}[Unsigned LEB128 encoding]
\label{def:leb128}
For a non-negative integer $x$, write $x$ in base $2^7$ as
$x = \sum_{i=0}^{\ell-1} g_i \cdot 2^{7i}$ with each \emph{group}
$g_i \in [0, 127]$, where $\ell$ is chosen minimally (i.e.\ $g_{\ell-1} \neq 0$,
except $\ell = 1$ and $g_0 = 0$ when $x = 0$). The encoding
$\mathrm{varint}(x)$ is the $\ell$-byte sequence $b_0 \| b_1 \| \cdots \|
b_{\ell-1}$ (\emph{little-endian group order}: the least-significant
7-bit group is emitted first) where
\[
  b_i = g_i \;\lor\; \bigl(\texttt{0x80} \text{ if } i < \ell-1 \text{ else } \texttt{0x00}\bigr),
\]
i.e.\ each byte carries 7 data bits in its low-order bits, and the
high-order (8th) bit is a \emph{continuation flag}: set to $1$ if another
group follows and $0$ on the final byte. Decoding reverses this: read
bytes, accumulating $\mathrm{group}(b_i) = b_i \,\&\, \texttt{0x7F}$
left-shifted by $7i$ into the result, until a byte with the continuation
flag clear is read.

\textbf{Canonicality.} A conforming decoder accepts \emph{only} the
canonical, minimal-length encoding: the unique $\ell$ satisfying
$2^{7(\ell-1)} \leq x < 2^{7\ell}$ for $x>0$ (equivalently, $b_{\ell-1} \neq
\texttt{0x00}$), or the single byte \texttt{0x00} for $x=0$. Decoding
\textbf{MUST reject}:
\begin{enumerate}
  \item \emph{Non-canonical (overlong) encodings} --- any encoding using
        more bytes than the minimal $\ell$ above, including a trailing
        group $g_{\ell-1} = 0$ appended after what would otherwise be the
        final byte;
  \item \emph{Inputs exceeding the maximum permitted integer width} ---
        this specification bounds every $\mathrm{varint}$-encoded value to
        $x < 2^{64}$ (consistent with the $\mathrm{be8}$ fixed-width field
        it replaces for block-mode tokens, Section~\ref{sec:wire-format-v7}),
        so a conforming decoder rejects any input requiring more than
        $\lceil 64/7 \rceil = 10$ groups, and truncated input (a
        continuation-flagged byte with no following byte).
\end{enumerate}
These rules make the mapping between integers and their encodings a
bijection on $[0, 2^{64})$, eliminating the ambiguity (multiple valid
encodings of the same integer) that overlong LEB128 encodings would
otherwise introduce.
\end{definition}

\subsection{Key}

A NAPQES key $\mathbf{k} = (k_1, k_2, \ldots, k_K)$ is an ordered tuple
of $K$ distinct primes drawn uniformly from $\mathcal{P}$ using a
cryptographically secure DRBG.  The serialised key material is:
\[
  \mathsf{key\_bytes} = \mathrm{be5}(k_1) \| \mathrm{be5}(k_2) \| \cdots \| \mathrm{be5}(k_K).
\]

The cardinality of $\mathcal{P}$ is approximately 586\,000; a 10-element
ordered tuple provides a key space of:
\[
  \frac{|\mathcal{P}|!}{(|\mathcal{P}|-10)!} \approx 1.1 \times 10^{58} \approx 2^{196}.
\]

The relevant security quantity for Theorem~\ref{thm:ind-cpa} below is not
the raw bit-length $8\cdot 5K = 40K$ of $\mathsf{key\_bytes}$, but the
\emph{min-entropy} of $\mathbf{k}$ as sampled by $\mathsf{KeyGen}$:
\[
  H_\infty(\mathbf{k}) \;=\; \log_2\!\left(\frac{|\mathcal{P}|!}{(|\mathcal{P}|-K)!}\right)
  \;\approx\; K\log_2|\mathcal{P}|,
\]
i.e.\ exactly $\log_2$ of the ordered-tuple key-space size above
($\approx 196$ bits for $K=10$, $\approx 19$ bits for $K=1$). Because
$\mathsf{kb}$ is a deterministic, injective encoding of $\mathbf{k}$,
$H_\infty(\mathsf{kb}) = H_\infty(\mathbf{k})$: injective post-processing
of a random variable never increases its min-entropy above the source's,
and here it does not decrease it either (Remark~\ref{rem:min-K} discusses
the one place this is not quite free --- HMAC's own key-length handling
for $5K>64$ bytes). This distinction --- $H_\infty(\mathbf{k}) \ll 40K$ ---
is the subject of audit finding \textbf{CVF8} below.

\subsection{Domain-Separated HMAC Derivation}

\begin{remark}[CVF2 fix: unified domain-first layout]
An earlier revision of this construction used a \emph{nonce-first} layout
$N \| d \| \mathsf{ctx}$ for domains $\{0,1,2,4,5,6,7\}$ but a
\emph{domain-first}, AAD-binding layout for the authentication-tag and
streaming-AE domains ($\{3,8,9\}$). Because Table~\ref{tab:domains} and this
section stated only the nonce-first formula, reading the wire-format tag
definition (Section~\ref{sec:wire-format-v7}) through that formula silently
dropped the AAD from the modelled HMAC input, and Remark~\ref{rem:domsep}
had to reason about the two layouts as separate groups, producing a
probabilistic cross-group collision term instead of an unconditional
separation argument. This was audit finding \textbf{CVF2}. The fix below
unifies all ten domains onto a single domain-first layout, after which
Remark~\ref{rem:domsep} needs only a single, unconditional injectivity
argument.
\end{remark}

All internal values are computed as:
\[
  \mathrm{Derive}_d(\mathsf{kb}, N, \mathsf{ctx}) = \mathrm{HMAC}(\mathsf{kb},\ d \| N \| \mathsf{ctx}),
\]
where $d \in \{0,1,\ldots,9\}$ is a 1-byte domain separator and every
variable-width field inside $\mathsf{ctx}$ is length-prefixed (e.g.
$\mathrm{be4}(|A|) \| A$ for AAD $A$).

\begin{table}[h]
\centering
\begin{tabular}{@{}clll@{}}
\toprule
Domain & Function & Context & Output \\
\midrule
\texttt{0x00} & Noise position & $\mathrm{be5}(ct\_pos)$ & Boolean \\
\texttt{0x01} & Real-token addend & $\mathrm{be5}(real\_idx)$ & $[1, k_i - 1]$ \\
\texttt{0x02} & Noise probability & (none) & $[0.75, 0.99]$ \\
\texttt{0x03} & Auth tag (block) & $\mathrm{be4}(\lvert aad \rvert) \| aad \| \widetilde{B}$ & 32 bytes \\
\texttt{0x04} & Noise character & $\mathrm{be5}(ct\_pos)$ & $[32, 127]$ \\
\texttt{0x05} & Noise-token addend & $\mathrm{be5}(ct\_pos)$ & $[1, k_i - 1]$ \\
\texttt{0x06} & Padding codepoint & $\mathrm{be4}(pad\_idx)$ & $[32, 126]$ \\
\texttt{0x07} & Varint keystream & $\mathrm{be4}(block\_idx)$ & 32-byte block \\
\texttt{0x08} & Chunk tag (stream AE) & $\mathrm{be4}(\lvert aad \rvert) \| aad \| \mathrm{be4}(i) \| \tilde{C}_i$ & 32 bytes \\
\texttt{0x09} & Final sentinel tag (stream AE) & $\mathrm{be4}(\lvert aad \rvert) \| aad \| \mathrm{be4}(n_{\mathrm{chunks}})$ & 32 bytes \\
\bottomrule
\end{tabular}
\caption{Domain-separated HMAC derivation schedule (CVF2 fix: unified domain-first layout, $\mathrm{Derive}_d(\mathsf{kb},N,\mathsf{ctx}) = \mathrm{HMAC}(\mathsf{kb}, d\|N\|\mathsf{ctx})$ for every domain; the nonce $N$ is applied uniformly by $\mathrm{Derive}_d$ and so is omitted from the Context column above). Domains \texttt{0x08}--\texttt{0x09} are used exclusively by the online-AE streaming format (Section~\ref{sec:streaming-ae}).}
\label{tab:domains}
\end{table}

\subsection{Noise Probability}

\[
  t = \frac{\mathrm{Derive}_{0x02}(\mathsf{kb}, N, \varepsilon)[0:8]_{\mathrm{be64}}}{2^{64}},
  \quad
  p = 0.75 + t \cdot 0.24.
\]

\subsection{Plaintext Padding}

Let $n = \lvert P \rvert$ (plaintext length in codepoints).  Define:
\[
  B = \max(16,\ 2^{\lceil \log_2(n+1) \rceil}), \qquad P_d = [n \gg 8,\ n \wedge 0\mathrm{FF}] \| P \| \mathbf{pad},
\]
where $\mathbf{pad} = (\mathbf{pad}_0, \ldots, \mathbf{pad}_{B-n-1})$ is a
$(B - n)$-codepoint sequence, each codepoint HMAC-derived via domain
\texttt{0x06} as made precise by Definition~\ref{def:padcodepoint} below.
The total padded length is $n + 2 + (B - n) = B + 2$ codepoints.

\begin{remark}[Audit finding CVF18: the padding-codepoint derivation formula was missing]
\label{rem:cvf18}
The text above previously named domain \texttt{0x06} and its output range
$[32,126]$ (Table~\ref{tab:domains}) but never stated how the 32-byte
output of $\mathrm{Derive}_{0x06}(\mathsf{kb}, N, \mathrm{be4}(pad\_idx))$
is reduced to a single codepoint in $[32,126]$ --- which bytes are taken
and what mapping is applied. Without that formula the padding is not
reproducible from the specification alone, and an unstated reduction into
a 95-value range is exactly the kind of construction that can introduce
modulo bias if done naively. Definition~\ref{def:padcodepoint} below
states the formula explicitly, matching the reference implementation
(\texttt{napqes.py}'s \texttt{\_pad\_message}), and Remark~\ref{rem:modbias}
is extended to cover its bias.
\end{remark}

\begin{definition}[Padding codepoint --- domain \texttt{0x06}]
\label{def:padcodepoint}
Let $w = \mathrm{Derive}_{0x06}(\mathsf{kb}, N, \mathrm{be4}(pad\_idx))[0{:}4]_{\mathrm{be32}} \in [0, 2^{32})$,
i.e.\ the first 4 bytes of the 32-byte HMAC output, interpreted as an
unsigned big-endian integer (the same $[0{:}n]_{\mathrm{be}k}$ convention
used by Definitions~\ref{def:noisechar}--\ref{def:realaddend}). Then
\[
  \mathbf{pad}_{pad\_idx} = (w \bmod 95) + 32 \ \in [32, 126].
\]
\end{definition}

\subsection{Token Construction}

\begin{remark}[Audit finding CVF6: token-emission pseudocode's helper primitives were undefined]
\label{rem:cvf6}
The token-emission loop below invokes four helper primitives ---
\texttt{is\_noise} (domain \texttt{0x00}), \texttt{derive\_noise\_char}
(\texttt{0x04}), \texttt{derive\_noise\_addend} (\texttt{0x05}), and
\texttt{derive\_real\_addend} (\texttt{0x01}) --- that were previously named
only by Table~\ref{tab:domains}'s context/range columns, with no formula
stating how the 32-byte $\mathrm{Derive}_d$ output is reduced to the
required Boolean decision or bounded integer. As written, the pseudocode
could not be executed and the token construction was not reproducible from
the paper alone --- and the exact reduction (and hence any modulo bias) was
unspecified for every one of these primitives. This was audit finding
\textbf{CVF6}. Definitions~\ref{def:isnoise}--\ref{def:realaddend} below
give each primitive as an explicit, byte-exact reduction of
$\mathrm{Derive}_d(\mathsf{kb},N,\mathsf{ctx})$, matching the reference
implementations exactly, and Remark~\ref{rem:modbias} bounds the resulting
modulo bias for each. This is a specification-completeness fix only --- no
algorithm or wire-format change; the reference implementations already
implemented exactly the formulas given below.
\end{remark}

For each padded codepoint $c$ (with current $ct\_pos$ and $real\_idx$),
the token is emitted by the loop below, which invokes the four helper
primitives \texttt{is\_noise}, \texttt{derive\_noise\_char},
\texttt{derive\_noise\_addend}, and \texttt{derive\_real\_addend}.
Definitions~\ref{def:isnoise}--\ref{def:realaddend} give each primitive
as an explicit reduction of $\mathrm{Derive}_d(\mathsf{kb},N,\mathsf{ctx})$
(Table~\ref{tab:domains}) into its target range, so that the loop is
fully specified and reproducible byte-for-byte.

Throughout, $[0{:}n]_{\mathrm{be}k}$ denotes the first $n$ bytes of a
32-byte HMAC-SHA256 digest, interpreted as an unsigned big-endian integer
in $[0, 2^{k})$ where $k = 8n$ (e.g.\ $[0{:}4]_{\mathrm{be32}} \in
[0,2^{32})$, matching the $[0{:}8]_{\mathrm{be64}}$ convention already used
for the noise probability $t$ above).

\begin{definition}[\texttt{is\_noise} --- domain \texttt{0x00}]
\label{def:isnoise}
Let $u = \mathrm{Derive}_{0x00}(\mathsf{kb}, N, \mathrm{be5}(ct\_pos))[0{:}8]_{\mathrm{be64}} / 2^{64} \in [0,1)$
and let $p$ be the noise probability of the previous subsection. Then
\[
  \texttt{is\_noise}(ct\_pos) = \mathbf{1}[\,u < p\,].
\]
This is a strict-inequality threshold test on a 64-bit fraction, not a
modular reduction, so it carries no modulo bias; the only deviation from a
continuous uniform draw is the $2^{-64}$ discretisation granularity, which
is cryptographically negligible.
\end{definition}

\begin{definition}[\texttt{derive\_noise\_char} --- domain \texttt{0x04}]
\label{def:noisechar}
Let $w = \mathrm{Derive}_{0x04}(\mathsf{kb}, N, \mathrm{be5}(ct\_pos))[0{:}4]_{\mathrm{be32}} \in [0,2^{32})$. Then
\[
  \texttt{derive\_noise\_char}(ct\_pos) = (w \bmod 96) + 32 \ \in [32,127].
\]
\end{definition}

\begin{definition}[\texttt{derive\_noise\_addend} --- domain \texttt{0x05}]
\label{def:noiseaddend}
Let $v = \mathrm{Derive}_{0x05}(\mathsf{kb}, N, \mathrm{be5}(ct\_pos))[0{:}4]_{\mathrm{be32}} \in [0,2^{32})$. Then, for key element $k$,
\[
  \texttt{derive\_noise\_addend}(ct\_pos, k) = (v \bmod (k-1)) + 1 \ \in [1, k-1].
\]
\end{definition}

\begin{definition}[\texttt{derive\_real\_addend} --- domain \texttt{0x01}]
\label{def:realaddend}
Let $v = \mathrm{Derive}_{0x01}(\mathsf{kb}, N, \mathrm{be5}(real\_idx))[0{:}4]_{\mathrm{be32}} \in [0,2^{32})$. Then, for key element $k$,
\[
  \texttt{derive\_real\_addend}(real\_idx, k) = (v \bmod (k-1)) + 1 \ \in [1, k-1].
\]
\end{definition}

\begin{remark}[Modulo bias of the finite-range reductions]
\label{rem:modbias}
Definitions~\ref{def:padcodepoint}, \ref{def:noisechar}--\ref{def:realaddend}
reduce a value $v$ drawn (pseudo-)uniformly from $[0,2^{32})$ into $[0,m)$ via
$v \bmod m$ before shifting into the target range. For any modulus
$m \le 2^{32}$, every residue class occurs either
$\lfloor 2^{32}/m \rfloor$ or $\lceil 2^{32}/m \rceil$ times among the
$2^{32}$ possible values of $v$, so each residue's probability deviates
from the ideal $1/m$ by at most $2^{-32}$ in absolute terms, i.e.\ by a
factor of at most $m/2^{32}$ relative to $1/m$. For
\texttt{derive\_noise\_char} ($m = 96$), this bound gives a relative
deviation below $3 \times 2^{-26} < 2^{-24}$. For the padding codepoint
(Definition~\ref{def:padcodepoint}, $m = 95$; audit finding CVF18,
Remark~\ref{rem:cvf18}), the same bound gives a relative deviation below
$95/2^{32} < 2^{-25}$. For \texttt{derive\_noise\_addend} and
\texttt{derive\_real\_addend}, $m = k - 1$ with $k \in \mathcal{P}
\subset [10^6, 9.9\times 10^6] \subset [0, 2^{24})$, giving a relative
deviation below $2^{24}/2^{32} = 2^{-8}$. In all cases the bias is
statistically dominated by, and computationally no more exploitable
than, the underlying PRF advantage against HMAC-SHA256, and is therefore
negligible for the security claims of Section~\ref{sec:security}; no
rejection sampling is required.
\end{remark}

\begin{lstlisting}[language=Python, caption={Token emission loop (pseudocode)}]
while is_noise(ct_pos):          # Definition 1, domain 0x00
    k        = key[real_idx % K]
    noise_c  = derive_noise_char(ct_pos)      # Definition 2, domain 0x04; in [32,127]
    noise_a  = derive_noise_addend(ct_pos, k) # Definition 3, domain 0x05; in [1,k-1]
    emit  noise_c * k + noise_a
    ct_pos += 1

k = key[real_idx % K]
a = derive_real_addend(real_idx, k)           # Definition 4, domain 0x01; in [1,k-1]
emit  c * k + a
ct_pos += 1;  real_idx += 1
\end{lstlisting}

\begin{remark}[Retracted: the GCD claim below was vacuous (audit finding CVF4)]
  An earlier revision claimed that, because $k$ is prime and
  $a \in [1, k-1]$, every real token satisfies $\gcd(\text{token}, k) < k$,
  and that the same holds for noise tokens, making real and noise tokens
  ``computationally indistinguishable under GCD analysis without the key.''
  This is vacuous: since $k$ is prime and every addend
  $a \in [1, k-1]$ by construction (Table~\ref{tab:domains}, domains
  \texttt{0x01}/\texttt{0x05}), $\gcd(a, k) = 1$ unconditionally, and
  $\gcd(\text{token}, k) = \gcd(c \cdot k + a,\ k) = \gcd(a, k) = 1$ for
  \emph{every} token --- real or noise, with or without the key. The
  statement requires no HMAC output and no noise-schedule design to hold,
  so it establishes no indistinguishability property and is removed. The
  layer's genuine, distinct contribution --- length-decorrelation from
  plaintext content, not GCD-level structural indistinguishability --- is
  stated and proved in Lemma~\ref{lem:tar} below.
\end{remark}

\subsection{Wire Format (Version 7 --- CVF1 fix)}
\label{sec:wire-format-v7}

\begin{remark}[Audit finding CVF1 and the fix below]
Version 6 of this construction encoded each token with
$\mathrm{varint}(\cdot)$ (LEB128). Because $\mathrm{token} = c \cdot k + a$,
the LEB128 byte-length of a token grows with the plaintext codepoint $c$: a
codepoint near 1 yields a $\approx 2^{20}$ token ($\approx 3$ LEB128 bytes),
while a codepoint near $2^{16}-1$ yields a $\approx 2^{36}$ token
($\approx 6$ LEB128 bytes). Two plaintexts with equal padded codepoint
count but different codepoint values therefore produced token blobs
$B_0, B_1$ of \emph{different byte-length} with overwhelming probability,
so the masked ciphertext $\widetilde{B}$ leaked plaintext content through
its length alone --- an IND-CPA distinguishing advantage $\approx 1$. This
invalidated the ``$B$ is the same length for both plaintexts'' step of the
hiding lemma (Lemma~\ref{lem:hiding} below). The fix, described here, is to
replace $\mathrm{varint}(\cdot)$ with a fixed-width encoding so that
$\lvert B \rvert$ is a function only of the token \emph{count}, which in
turn is a function of the padded codepoint count and the HMAC-derived,
content-independent noise schedule (domain \texttt{0x00}) --- never of
codepoint values. Lemma~\ref{lem:hiding} is restated and re-proved below
under this v7 encoding; Section~\ref{sec:hiding-lemma-erratum} discusses the
v6 argument's error precisely.
\end{remark}

Each token is serialised as $\mathrm{be8}(\mathrm{token})$ instead of
$\mathrm{varint}(\mathrm{token})$:
\[
  B = \mathrm{be8}(\mathrm{token}[0]) \| \mathrm{be8}(\mathrm{token}[1]) \| \cdots \| \mathrm{be8}(\mathrm{token}[N-1]),
\]
so $\lvert B \rvert = 8N$ where $N$ (real $+$ noise tokens) depends only on
the padded codepoint count and $(\mathsf{kb}, N_{\text{once}})$ via the
noise-position oracle (domain \texttt{0x00}) --- not on codepoint values.
The overall ciphertext is unchanged in shape:
\[
  C = N \| \widetilde{B} \| T,
\]
where $N$ is a 16-byte random nonce and $\widetilde{B}$ is the
domain-\texttt{0x07} XOR-masked fixed-width token blob, defined explicitly
as
\[
  \widetilde{B} \;=\; B \,\oplus\, \mathrm{ks}_{0x07}(N)\bigl[0 : \lvert B \rvert\bigr],
  \qquad
  \mathrm{ks}_{0x07}(N) \;=\; \mathrm{Derive}_{0x07}(\mathsf{kb}, N, \mathrm{be4}(0)) \,\|\,
  \mathrm{Derive}_{0x07}(\mathsf{kb}, N, \mathrm{be4}(1)) \,\|\, \cdots,
\]
i.e.\ the encode-then-mask order is: (1) assemble the fixed-width token
blob $B$ from the token vector $(\mathrm{token}[0], \ldots,
\mathrm{token}[N-1])$ produced by the token-construction loop above
(Definitions~\ref{def:isnoise}--\ref{def:realaddend}); (2) generate the
domain-\texttt{0x07} keystream $\mathrm{ks}_{0x07}(N)$ as the
concatenation of 32-byte $\mathrm{Derive}_{0x07}$ blocks indexed
$0, 1, 2, \ldots$ (Table~\ref{tab:domains}), generating only as many
blocks as needed to cover $\lvert B \rvert$ bytes and truncating the final
block to the remaining length; (3) XOR the two byte strings of equal
length $\lvert B \rvert$. This equation is proved to yield a value
statistically independent of, and (for fixed $N$) uniformly distributed
given, $B$ in the ``Independence of $\mathrm{ks}(N^*)$ from $B$'' argument
of Lemma~\ref{lem:hiding}'s proof, which this definition forward-references.
\begin{remark}[Audit finding CVF19: $\widetilde{B}$ used without a defining equation at its point of introduction]
\label{rem:cvf19}
This subsection previously introduced $\widetilde{B}$ only in prose ---
``the domain-\texttt{0x07} XOR-masked fixed-width token blob'' --- with no
defining equation at this point; the relation
$\widetilde{B} = B \oplus \mathrm{ks}_{0x07}(N)$, the encode-then-mask
operation order, and the keystream length/truncation rule were stated
only later, inside Lemma~\ref{lem:hiding}'s proof
(Section~\ref{sec:hiding-lemma-erratum} and surrounding text). As a
result, $C = N \| \widetilde{B} \| T$ could not be reproduced from this
section alone. The equation above closes that gap by stating the
definition explicitly here, with forward references to the token
construction (Definitions~\ref{def:isnoise}--\ref{def:realaddend}) and the
domain-\texttt{0x07} keystream construction it depends on.
\end{remark}
\[
  T = \mathrm{Derive}_{0x03}\bigl(\mathsf{kb},\ N,\ \mathrm{be4}(|A|) \| A \| \widetilde{B}\bigr)
\]
is the 32-byte HMAC-SHA256 authentication tag, with $A$ the (possibly empty)
AAD (CVF2 fix: this tag definition now matches Table~\ref{tab:domains}
exactly and explicitly binds $A$, whereas an earlier revision's inline
formula, $T = \mathrm{Derive}_{0x03}(\mathsf{kb}, N \| \widetilde{B})$, omitted
the AAD and was — read through the then-stated nonce-first
$\mathrm{Derive}_d$ formula — inconsistent with the AAD-binding tag the
reference code actually computed). The minimum valid ciphertext length is
48 bytes.

\subsection{Online-AE Streaming Format (v6s-ae)}
\label{sec:streaming-ae}

The basic streaming format (\texttt{encrypt\_stream}) carries a single tag at the
end of the stream, causing \texttt{decrypt\_stream} to release plaintext before
the tag is verified (RUP, formerly CAV-001).  The online-AE streaming format
eliminates this by attaching a per-chunk HMAC-SHA256 tag to each segment of the
masked varint blob.

\[
  S = N \;\|\; \Bigl[\,\mathrm{be4}(\ell_i) \| \widetilde{C}_i \| T_i\,\Bigr]_{i=0}^{n-1} \;\|\; \mathrm{be4}(0) \| T_{\mathrm{final}},
\]
where each $\widetilde{C}_i$ is at most 1024 bytes of the masked varint blob and (CVF2
fix: nonce moved to the fixed byte $1..16$ offset shared with every other domain):
\begin{align*}
  T_i &= \mathrm{HMAC}\!\left(\mathsf{kb},\ \texttt{0x08}
          \| N \| \mathrm{be4}(|A|) \| A \| \mathrm{be4}(i) \| \widetilde{C}_i\right), \\
  T_{\mathrm{final}} &= \mathrm{HMAC}\!\left(\mathsf{kb},\ \texttt{0x09}
          \| N \| \mathrm{be4}(|A|) \| A \| \mathrm{be4}(n)\right).
\end{align*}
The sentinel frame has $\ell = 0$; it signals end-of-stream.
$\texttt{decrypt\_stream\_ae}$ verifies $T_i$ before yielding any plaintext from chunk $i$.
$T_{\mathrm{final}}$ binds the total chunk count $n$, preventing silent truncation at a
chunk boundary.  Chunk index $i$ in each $T_i$ prevents chunk-reordering attacks.

\subsection{Format Applicability and Normative Status (CVF7 fix)}
\label{sec:format-applicability}

\begin{remark}[Audit finding CVF7]
The preceding subsections define three distinct wire encodings — block-mode
Wire Format Version 7 ($C = N \| \widetilde{B} \| T$,
Section~\ref{sec:wire-format-v7}), the basic streaming format
($\texttt{encrypt\_stream}$/$\texttt{decrypt\_stream}$, a single trailing
tag), and the online-AE streaming format (v6s-ae,
Section~\ref{sec:streaming-ae}, per-chunk tags) — without ever stating
their relationship: whether they are interchangeable options, whether
block-mode and streaming coexist by design, or how a decryptor determines
which of the three a given ciphertext uses, since none of the three wire
layouts carries a version/format discriminator byte. This was audit finding
\textbf{CVF7}. Because the basic streaming format is RUP-prone (CAV-001,
Section~\ref{sec:caveats}), leaving the relationship unstated has a security
consequence: absent a stated selection rule, a receiver could be induced to
invoke the unsafe decoder on a ciphertext produced by (or claimed to come
from) the safe encoder. The remainder of this subsection states the rule
that was missing.
\end{remark}

\paragraph{Block mode and streaming mode coexist by design; both are
normative.} They are not alternative encodings of one message class chosen
by preference — they serve two disjoint use cases:
\begin{itemize}
  \item \textbf{Block mode} (Wire Format v7, Section~\ref{sec:wire-format-v7})
        is normative for bounded messages that fit in memory and require
        length-bucket obfuscation via padding. It is the sole format
        underlying $\mathsf{Enc}/\mathsf{Dec}$
        (Definition~\ref{def:aead-triple}) and the sole format covered by
        the IND-CPA (Theorem~\ref{thm:ind-cpa}) and INT-CTXT theorems of
        Section~\ref{sec:security}.
  \item \textbf{Streaming mode} is normative for unbounded or
        memory-constrained plaintext that cannot be buffered whole. It
        never applies block padding, so a streaming ciphertext always
        discloses the exact plaintext length and is never interoperable
        with a block-mode ciphertext (different framing, different token
        encoding; Section~\ref{sec:streaming-ae}).
\end{itemize}

\paragraph{Within streaming mode, only the online-AE format is normative;
the basic streaming format is deprecated and forbidden for new
ciphertexts.}
\begin{itemize}
  \item \textbf{The online-AE streaming format (v6s-ae,
        Section~\ref{sec:streaming-ae}) is, as of this fix, the sole
        normative streaming format.} It MUST be used for every new
        streaming deployment.
  \item \textbf{The basic streaming format
        ($\texttt{encrypt\_stream}$/$\texttt{decrypt\_stream}$) is
        superseded, deprecated, and its use is forbidden for producing new
        ciphertext.} It is retained solely so that streams produced before
        this fix remain decryptable; no protocol may be designed to accept
        it for new traffic, and no new ciphertext may be issued in this
        format. This upgrades CAV-001 (Section~\ref{sec:caveats}) from
        ``fixed, with the RUP-prone format still permitted at the caller's
        discretion'' to ``fixed by supersession: the RUP-prone format is
        deprecated and forbidden going forward.''
\end{itemize}

\paragraph{Format selection is an out-of-band API contract, not an in-band
discriminator.} None of the three wire layouts carries a version/format
byte, and this fix does not add one. Exactly as with key and nonce
agreement, the two endpoints of a NAPQES deployment agree out-of-band, at
the protocol/call-site level, on which of block mode,
$\texttt{encrypt\_stream\_ae}$/$\texttt{decrypt\_stream\_ae}$, or the
deprecated $\texttt{encrypt\_stream}$/$\texttt{decrypt\_stream}$ is in use
for a given channel or message class — the same way a caller of any AEAD
library chooses between a one-shot and a streaming construction rather than
having the library sniff the format from ciphertext bytes. This is a
deliberate design choice and not an omission: a self-describing in-band
discriminator would itself have to be authenticated before a decryptor
could safely act on it, which reintroduces the same
trust-before-verification problem RUP already poses. Three consequences
follow: (1) a decryptor is never expected to auto-detect format from
ciphertext content alone; (2) a protocol built on NAPQES MUST NOT implement
a ``try v6s-ae, then fall back to the basic streaming format'' negotiation,
since that downgrade path would recreate exactly the induced-unsafe-decode
hazard this fix closes; and (3) any future in-band discriminator, should
one be added, would itself require authentication under the tag before a
decryptor may act on it — that mechanism is deferred as out of scope for
this fix.

\subsection{The Legacy v7 AEAD Algorithm Triple (Backward Compatibility)}
\label{sec:algorithm-triple}

\begin{remark}[Audit finding CVF5]
The subsections above state notation, key generation, per-domain HMAC
derivation, noise probability, padding, the token-emission loop, and the
wire format, but nowhere assembled them into the formal algorithm triple
that an AEAD definition requires (cf.\ the Ascon v1.2
submission~\cite{ascon2021}), and
Section~\ref{sec:security} already invokes $\mathrm{NAPQES}.\mathsf{Enc}$
and $\mathrm{NAPQES}.\mathsf{Dec}$ (e.g.\ the INT-CTXT proof,
Section~\ref{sec:security}) without those symbols ever having been formally
defined — a key-generation/setup algorithm, an encryption algorithm
$\mathsf{Enc}(K,N,A,M)\to C$, and a decryption algorithm
$\mathsf{Dec}(K,N,A,C)\to M$ or $\bot$, each with explicitly typed
inputs/outputs and a stated correctness condition, were left implicit and
scattered across the prose and pseudocode above. This was audit finding
\textbf{CVF5}. Definition~\ref{def:aead-triple} below assembles the pieces
already introduced into that triple; no algorithm, encoding, or derivation
changes — this is a specification-assembly fix, not a code or wire-format
change.
\end{remark}

\begin{definition}[NAPQES AEAD algorithm triple]\label{def:aead-triple}
NAPQES is the triple $(\mathsf{KeyGen}, \mathsf{Enc}, \mathsf{Dec})$ over the
following typed spaces:
\begin{itemize}
  \item Key space $\mathcal{K}$: ordered $K$-tuples of distinct elements of
        $\mathcal{P}$ (Section~\ref{sec:construction}, ``Key'').
  \item Nonce space $\mathcal{N} = \{0,1\}^{128}$.
  \item AAD space $\mathcal{A} = \{0,1\}^{*}$.
  \item Message space $\mathcal{M}$: finite sequences of Unicode codepoints
        (Python \texttt{str} values, read via $\mathrm{ord}(\cdot)$) of length
        at most \texttt{MAX\_PLAINTEXT\_CODEPOINTS} $= 2^{16}-1$, the limit
        imposed by the 2-codepoint big-endian length prefix written during
        padding (Section~\ref{sec:construction}, ``Plaintext Padding'').
  \item Ciphertext space $\mathcal{C} = \{0,1\}^{*}$; every output of
        $\mathsf{Enc}$ has length $16 + 8N + 32$ bytes, where $N$ is the
        real-plus-noise token count for that message (minimum 48 bytes).
\end{itemize}

\paragraph{$\mathsf{KeyGen}(1^{\lambda}) \to \mathbf{k}$.}
Draw $K$ distinct primes $k_1,\ldots,k_K$ uniformly without replacement from
$\mathcal{P}$ by CSPRNG-driven rejection sampling
(\texttt{generate\_prime\_numbers}); output
$\mathbf{k} = (k_1,\ldots,k_K) \in \mathcal{K}$ and set
$\mathsf{kb} = \mathsf{key\_bytes}(\mathbf{k})$ as in
Section~\ref{sec:construction}, ``Key''.

\paragraph{$\mathsf{Enc}(\mathbf{k}, N, A, M) \to C$.}
On input $\mathbf{k}\in\mathcal{K}$, $N\in\mathcal{N}$, $A\in\mathcal{A}$,
$M\in\mathcal{M}$: (1) compute $\mathsf{kb}$; (2) pad $M$ to $P_d$
(``Plaintext Padding''); (3) run the token-emission loop over $P_d$ using
$(\mathsf{kb}, N)$ (``Token Construction'') to obtain the token sequence
$\mathrm{token}[0..N{-}1]$; (4) serialise
$B = \mathrm{be8}(\mathrm{token}[0])\|\cdots\|\mathrm{be8}(\mathrm{token}[N{-}1])$
(``Wire Format, Version 7''); (5) mask
$\widetilde{B} = B \oplus \mathrm{Derive}_{0x07}(\mathsf{kb}, N, \cdot)$;
(6) compute the tag
$T = \mathrm{Derive}_{0x03}(\mathsf{kb}, N, \mathrm{be4}(|A|)\|A\|\widetilde{B})$;
(7) output $C = N \| \widetilde{B} \| T$.

This is the algorithm the audit's recommended signature,
$\mathsf{Enc}(K,N,A,M)\to C$, refers to, and it is exactly the internal,
explicit-nonce encryption routine used by the FIPS power-on self-test and
the KAT harness. The public, randomized API,
$\mathsf{Encrypt}(\mathbf{k}, A, M) \triangleq \mathsf{Enc}(\mathbf{k}, N, A, M)$
for $N \xleftarrow{\$} \{0,1\}^{128}$ sampled internally by a CSPRNG
(\texttt{encrypt\_bytes}/\texttt{encrypt\_str}), is a randomized wrapper
around $\mathsf{Enc}$ that never exposes $N$ as a caller-chosen input
(the explicit-nonce $\mathsf{Enc}$ itself is not part of the public API —
see \textbf{CVF3} for why caller-chosen nonces were restricted). This
distinction — a deterministic, nonce-keyed core algorithm versus a
randomized public wrapper that samples the nonce — is why nonce reuse
(CVF3) is a DRBG-failure/replay hazard rather than a caller-misuse hazard
in normal use, and why Theorem~\ref{thm:ind-cpa}'s experiment (which
queries the public, randomized encryption oracle) and Lemma~\ref{lem:tar}'s
per-nonce framing both apply to the same underlying $\mathsf{Enc}$.

\paragraph{$\mathsf{Dec}(\mathbf{k}, A, C) \to M \text{ or } \bot$.}
On input $\mathbf{k}\in\mathcal{K}$, $A\in\mathcal{A}$,
$C\in\{0,1\}^{*}$: (1) if $|C| < 48$, output $\bot$; (2) parse
$C = N \| \widetilde{B} \| T'$ with $N \in \{0,1\}^{128}$ the leading 16
bytes and $T'$ the trailing 32 bytes; (3) compute $\mathsf{kb}$ and
$T = \mathrm{Derive}_{0x03}(\mathsf{kb}, N, \mathrm{be4}(|A|)\|A\|\widetilde{B})$;
(4) if $T \neq T'$ (compared in constant time), output $\bot$; (5) else
unmask $B = \widetilde{B} \oplus \mathrm{Derive}_{0x07}(\mathsf{kb}, N, \cdot)$,
deserialise the fixed-width tokens, invert the token-emission loop using the
same $(\mathsf{kb}, N)$-derived noise-position oracle and addends to recover
$P_d$, strip the length prefix and padding, and output $M$.

We write $\mathsf{Dec}$ as taking $(\mathbf{k}, A, C)$ rather than the
audit's suggested $(\mathbf{k}, N, A, C)$ because NAPQES's wire format
embeds $N$ as $C$'s first 16 bytes (step 2 above) rather than transmitting
it out-of-band; this matches the actual \texttt{decrypt\_bytes(ciphertext,
key, aad)} call signature and the $\mathsf{Dec}(\mathbf{k}, c^{*}, aad^{*})$
notation already used in the INT-CTXT proof (Section~\ref{sec:security}),
and is a trivial reparameterisation of the audit's four-argument form under
$N := C[0{:}16]$ — no security-relevant difference results from embedding
$N$ in $C$ versus passing it separately, since both are fully determined by
the ciphertext, and the domain-separation argument of
Remark~\ref{rem:domsep} does not depend on which of the two calling
conventions is used.
\end{definition}

\begin{definition}[Correctness]\label{def:correctness}
For every $\lambda$, every $\mathbf{k} \gets \mathsf{KeyGen}(1^{\lambda})$,
every $N\in\mathcal{N}$, $A\in\mathcal{A}$, $M\in\mathcal{M}$:
\[
  \mathsf{Dec}\bigl(\mathbf{k},\, A,\, \mathsf{Enc}(\mathbf{k}, N, A, M)\bigr) = M.
\]
This holds by construction: $\mathsf{Enc}$'s token-emission loop (step 3)
and $\mathsf{Dec}$'s loop-inversion step (step 5) both compute the
noise-position oracle (domain \texttt{0x00}), the real/noise addends
(domains \texttt{0x01}/\texttt{0x05}), and the padding codepoints (domain
\texttt{0x06}) as pure, deterministic functions of $(\mathsf{kb}, N,
\mathrm{position})$; given the same $\mathbf{k}$ and $N$, both algorithms
therefore compute identical intermediate values at every position, so
token emission and token inversion are exact inverses whenever $T = T'$.
Correctness is additionally checked empirically by the cross-language KAT
corpus (\texttt{tests/kat/v6\_vectors.json}) and the full regression suite
across the Python, Rust, and C implementations
(Section~\ref{sec:implementation}).
\end{definition}

\begin{remark}[Authentication-failure output $\bot$]
The formal $\bot$ output of $\mathsf{Dec}$ is realised in all three
reference implementations as a raised exception
(\texttt{ValueError} in Python/Rust/C-equivalent error codes) rather than a
sentinel return value; SPEC.md \S10 (``Error model'') enumerates every
condition under which $\mathsf{Dec}$ raises rather than returns $M$,
including tag mismatch, undersized ciphertext, and truncated streaming
input. This is a standard, semantically equivalent realisation of $\bot$
(cf.\ most software AEAD APIs, e.g.\ \texttt{cryptography.hazmat}), and every
occurrence of ``$\mathsf{Dec}(\cdot)=\bot$'' or
``$\mathsf{Dec}(\cdot)\neq\bot$'' in Section~\ref{sec:security} below
should be read as, respectively, ``raises'' or ``does not raise''.
\end{remark}

\subsection{The NAPQES Construction, Version 2 (Default): Independent Key Schedule and Synthetic Nonce (CVF3 / CVF8 / CVF13 fix)}
\label{sec:v8}

\begin{remark}[Why v7 cannot close CVF3/CVF8/CVF13 without a key-schedule change]
\label{rem:v8-motivation}
Three audit findings trace back to a single design choice shared by every
wire format up to and including v7: \emph{one} secret, $\mathbf{k}$ (via
$\mathsf{kb}=\mathsf{key\_bytes}(\mathbf{k})$), plays two unrelated roles
--- it keys every domain-derivation HMAC call, \emph{and} it is the
arithmetic-layer prime tuple used directly by the token map $c\mapsto
c\cdot k+a$.
\begin{itemize}
  \item \textbf{CVF3.} Because every derived value is a deterministic
        function of $(\mathsf{kb}, N)$, a repeated nonce reproduces an
        identical keystream and addends; combined with the exact affine
        token map, this permits exact key recovery from two known
        plaintexts under one reused nonce (footnote to
        Table~\ref{tab:comparison}). A CSPRNG-sampled nonce only makes
        \emph{accidental} reuse improbable; it does nothing against DRBG
        failure, restart/replay, or process-snapshot reuse.
  \item \textbf{CVF8.} $\mathsf{kb}$ has only $H_\infty(\mathbf{k})\approx
        19.16K$ bits of min-entropy (Remark~\ref{rem:cvf8}), a non-standard
        HMAC-key hypothesis distinct from the textbook uniform-key PRF
        conjecture.
  \item \textbf{CVF13.} A reduction simulating $\mathsf{Enc}$ must know
        $\mathbf{k}$ to run the arithmetic layer, yet the same $\mathbf{k}$
        (via $\mathsf{kb}$) is supposed to be the PRF oracle's hidden key
        --- a reduction that already knows the ``hidden'' key cannot
        meaningfully distinguish it from random, so the PRF hop is vacuous
        as written (Remark~\ref{rem:cvf13}).
\end{itemize}
Remark~\ref{rem:cvf8-kdf-erratum} shows that deriving
$\mathsf{sk}=\mathrm{HKDF}(\mathsf{kb})$ --- the fix originally sketched
under CVF8 --- does \emph{not} resolve CVF13, because $\mathsf{sk}$ remains
a deterministic function of $\mathbf{k}$: knowing $\mathbf{k}$ still means
knowing $\mathsf{sk}$. All three findings are closed simultaneously only by
making the HMAC-keying secret and the arithmetic-layer secret
\emph{independently sampled} random variables, with no deterministic
function relating one to the other. This subsection specifies that
construction (v8) precisely.
\end{remark}

\begin{definition}[V8 key generation]\label{def:keygen-v8}
$\mathsf{KeyGen}_{v8}(1^\lambda) \to (\mathbf{k}, \mathsf{sk})$: draw
$\mathbf{k} = (k_1,\ldots,k_K)$ exactly as in $\mathsf{KeyGen}$
(Definition~\ref{def:aead-triple}), and, \emph{independently}, draw
$\mathsf{sk} \xleftarrow{\$} \{0,1\}^{256}$ uniformly via CSPRNG
(\texttt{rand::thread\_rng().fill\_bytes}, \texttt{generate\_v8\_key} in
\texttt{rust/src/lib.rs}). No function relates $\mathsf{sk}$ to
$\mathbf{k}$ or to $\mathsf{kb}=\mathsf{key\_bytes}(\mathbf{k})$; both
$\mathbf{k}$ and $\mathsf{sk}$ are secret key material that MUST be
generated, stored, and transmitted together.
\end{definition}

\begin{definition}[Synthetic nonce --- domain \texttt{0x0A}]\label{def:synthnonce}
For $\mathsf{sk}\in\{0,1\}^{256}$, $A\in\{0,1\}^{*}$, $M\in\mathcal{M}$
(encoded as UTF-8 bytes),
\[
  N_{v8}(\mathsf{sk}, A, M) \;=\; \mathrm{HMAC}\bigl(\mathsf{sk},\;
    \texttt{0x0A} \,\|\, \mathrm{be4}(|A|) \,\|\, A \,\|\, M\bigr)[0{:}16].
\]
Unlike domains \texttt{0x00}--\texttt{0x09}, this derivation takes no
nonce as input --- it is the step that \emph{produces} the nonce, in the
manner of RFC~5297 SIV / AES-GCM-SIV's synthetic IV, and is keyed only by
$\mathsf{sk}$, never by $\mathbf{k}$.
\end{definition}

\begin{definition}[V8 encryption and decryption]\label{def:aead-triple-v8}
$\mathsf{Enc}_{v8}(\mathbf{k}, \mathsf{sk}, A, M) \to C$: (1) compute
$N = N_{v8}(\mathsf{sk}, A, M)$ (Definition~\ref{def:synthnonce}); (2)
run $\mathsf{Enc}(\mathbf{k}, N, A, M)$ (Definition~\ref{def:aead-triple})
verbatim, with every occurrence of $\mathsf{kb}$ in that algorithm's steps
(2)--(6) replaced by $\mathsf{sk}$ --- i.e.\ padding (domain
\texttt{0x06}), the noise-position oracle (\texttt{0x00}), addends
(\texttt{0x01}/\texttt{0x05}), noise characters (\texttt{0x04}), the
\texttt{0x07} keystream, and the \texttt{0x03} tag are all keyed by
$\mathsf{sk}$; only the token map $c\mapsto c\cdot k+a$ itself still uses
$\mathbf{k}$ directly, exactly as in $\mathsf{Enc}$. Output $C$.
$\mathsf{Dec}_{v8}(\mathbf{k}, \mathsf{sk}, A, C) \to M \text{ or } \bot$:
run $\mathsf{Dec}(\mathbf{k}, A, C)$ verbatim with $\mathsf{kb}$ replaced
by $\mathsf{sk}$ throughout. Correctness
(Definition~\ref{def:correctness}) carries over unchanged, since the
substitution is uniform and the same $\mathsf{sk}$ is used by both
algorithms. Implemented as \texttt{encrypt\_bytes\_v8}/
\texttt{decrypt\_bytes\_v8} in \texttt{rust/src/lib.rs}; the wire-format
byte shape ($N\Vert\widetilde{B}\Vert T$) is unchanged, but v7 and v8
ciphertexts are \emph{not} interoperable with each other (different
keying and nonce derivation), and --- per the format-selection philosophy
of Section~\ref{sec:format-applicability} --- callers must agree
out-of-band on which schedule a given key/ciphertext pair uses.
\end{definition}

\begin{lemma}[V8 closes CVF3: nonce reuse across distinct messages is negligible]
\label{lem:v8-cvf3}
For any two queries $(A_1,M_1) \neq (A_2,M_2)$ to $\mathsf{Enc}_{v8}$
under the same $\mathsf{sk}$, $\Pr[N_{v8}(\mathsf{sk},A_1,M_1) =
N_{v8}(\mathsf{sk},A_2,M_2)]$ is bounded by
$\mathrm{Adv}^{\mathrm{PRF}}_{\mathrm{HMAC\text{-}SHA256}}(\mathcal{B}) +
2^{-128}$ (a PRF hop identical in form to Lemma~\ref{lem:hiding}'s hop,
plus the birthday bound on the truncated 128-bit output), for a PRF
adversary $\mathcal{B}$ with access to $\mathsf{sk}$ only via domain
\texttt{0x0A} queries. Consequently the affine-cancellation key-recovery
route of the Table~\ref{tab:comparison} footnote --- which requires two
\emph{distinct} known plaintexts under one \emph{reused} nonce --- applies
to $\mathsf{Enc}_{v8}$ only with this same negligible probability, rather
than with the $\approx 2^{-64}$ birthday probability of v7's independent
CSPRNG nonce. Re-querying $\mathsf{Enc}_{v8}$ on an \emph{identical}
$(A,M)$ pair reproduces the identical $N$ (and hence the identical $C$)
deterministically by Definition~\ref{def:synthnonce} --- this discloses
only plaintext equality, the standard, disclosed MRAE trade-off (cf.\
AES-GCM-SIV), not a key-recovery or confidentiality break.
\end{lemma}

\begin{lemma}[V8 closes CVF8: the standard uniform-key PRF assumption applies]
\label{lem:v8-cvf8}
By Definition~\ref{def:keygen-v8}, $\mathsf{sk}\xleftarrow{\$}\{0,1\}^{256}$
exactly, so $H_\infty(\mathsf{sk})=256$ and
$\mathrm{Adv}^{\mathrm{PRF}}_{\mathrm{HMAC\text{-}SHA256}}(\mathcal{B})$
against $\mathsf{sk}$ is the conventional, uniform-key HMAC-SHA256 PRF
advantage the literature states --- no key-guessing term
$q_F\cdot2^{-H_\infty(\cdot)}$ (Remark~\ref{rem:cvf8}) and no
PRF-under-non-uniform-$D$ caveat (Remark~\ref{rem:prf-d}) is needed for a
v8-restated Theorem~\ref{thm:ind-cpa}: simply replace $\mathsf{kb}$ by
$\mathsf{sk}$ throughout the proof of Theorem~\ref{thm:ind-cpa} and drop
the $q_F\cdot2^{-H_\infty(\mathbf{k})}$ term, since $\mathbf{k}$ no longer
keys any HMAC call the adversary can query.
\end{lemma}

\begin{lemma}[V8 closes CVF13: the reductions are genuinely simulatable]
\label{lem:v8-cvf13}
Let $\mathcal{O}$ be a PRF challenger's oracle, keyed either by a uniform
$\mathsf{sk}^*\xleftarrow{\$}\{0,1\}^{256}$ (real world) or by a truly
random function (ideal world). Construct $\mathcal{B}$ as follows:
$\mathcal{B}$ samples its own $\mathbf{k}'$ via
$\mathsf{KeyGen}_{v8}$'s prime-sampling step (Definition~\ref{def:keygen-v8}),
independently of $\mathcal{O}$; runs $\mathsf{Enc}_{v8}(\mathbf{k}',
\mathsf{sk}^*, A, M)$ for each of $\mathcal{A}$'s queries by computing the
arithmetic layer ($c\cdot k'_j+a$) itself using $\mathbf{k}'$, and
forwarding every domain-\texttt{0x00}--\texttt{0x0A} derivation call to
$\mathcal{O}$ in place of directly computing $\mathrm{HMAC}(\mathsf{sk}^*,\cdot)$.
Because $\mathsf{sk}^*$ is sampled independently of $\mathbf{k}'$
(Definition~\ref{def:keygen-v8}), $\mathcal{B}$ never needs to know
$\mathsf{sk}^*$ to run this simulation --- unlike the v7 case
(Remark~\ref{rem:cvf13}) and unlike the rejected
$\mathsf{sk}=\mathrm{HKDF}(\mathsf{kb})$ repair
(Remark~\ref{rem:cvf8-kdf-erratum}), where knowing the arithmetic key
implied knowing the HMAC key. When $\mathcal{O}$ is real,
$\mathcal{B}$'s simulation is distributed identically to
$\mathsf{Enc}_{v8}(\mathbf{k}, \mathsf{sk}^*, \cdot, \cdot)$ for the true,
independently-sampled $\mathbf{k}$ (both are: a uniformly-drawn prime
tuple, together with $\mathsf{Enc}_{v8}$ run under the real $\mathsf{sk}^*$
--- Lemma~\ref{lem:hiding}/Lemma~\ref{lem:tar} already establish that the
adversary's view depends on $\mathbf{k}$ only through the
$\mathsf{sk}^*$-masked blob and the content-independent token \emph{count},
not through $\mathbf{k}$'s specific value); when $\mathcal{O}$ is ideal,
it is distributed identically to the $G_1$ hop of
Theorem~\ref{thm:ind-cpa}/Theorem~\ref{thm:int-ctxt}. Hence
$|\Pr[\mathcal{A}\text{ wins real}]-\Pr[\mathcal{A}\text{ wins ideal}]| \leq
\mathrm{Adv}^{\mathrm{PRF}}_{\mathrm{HMAC\text{-}SHA256}}(\mathcal{B})$,
exactly the PRF hop Theorem~\ref{thm:ind-cpa} and
Theorem~\ref{thm:int-ctxt} require --- the simulation gap Remark~\ref{rem:cvf13}
identified for v7 does not arise for v8.
\end{lemma}

\begin{theorem}[IND-CPA, v8]\label{thm:ind-cpa-v8}
Let $\Pi_{v8}=(\mathsf{KeyGen}_{v8},\mathsf{Enc}_{v8},\mathsf{Dec}_{v8})$
(Definitions~\ref{def:keygen-v8}, \ref{def:synthnonce},
\ref{def:aead-triple-v8}) be the v8 NAPQES construction, with
$\mathsf{Encrypt}_{v8}(\mathbf{k},\mathsf{sk},A,M)\triangleq
\mathsf{Enc}_{v8}(\mathbf{k},\mathsf{sk},A,M)$ its (already randomized,
via the synthetic nonce) public encryption oracle, and let
$\mathbf{Exp}^{\mathrm{ind\text{-}cpa}}_{\Pi_{v8},\mathcal{A}}(\lambda)$ be
Definition~\ref{def:ind-cpa}'s experiment with $\mathsf{Encrypt}$ replaced
by $\mathsf{Encrypt}_{v8}$ throughout. For any PPT adversary $\mathcal{A}$
making at most $q$ encryption queries, there exists a PRF adversary
$\mathcal{B}_1'$ with similar running time such that
\[
  \mathrm{Adv}^{\mathrm{IND\text{-}CPA}}_{\Pi_{v8}}(\mathcal{A})
  \;\leq\;
  \mathrm{Adv}^{\mathrm{PRF}}_{\mathrm{HMAC\text{-}SHA256}}(\mathcal{B}_1')
  + \frac{q^2}{2^{128}},
\]
where $\mathrm{Adv}^{\mathrm{PRF}}_{\mathrm{HMAC\text{-}SHA256}}(\mathcal{B}_1')$
is the \emph{standard, uniform-key} HMAC-SHA256 PRF advantage --- with
\textbf{no} key-guessing term $q_F\cdot2^{-H_\infty(\mathbf{k})}$ and no
PRF-under-non-uniform-$D$ caveat.
\end{theorem}

\begin{proof}
Identical to the proof of Theorem~\ref{thm:ind-cpa}, with every occurrence
of $\mathsf{kb}$ replaced by $\mathsf{sk}$, subject to two changes
justified by Lemmas~\ref{lem:v8-cvf3}--\ref{lem:v8-cvf13} above. First,
Lemma~\ref{lem:key-guess}'s guessing event $\mathrm{Guess}$ (an adversary
offline-guessing the HMAC key) no longer applies: by
Lemma~\ref{lem:v8-cvf8}, $\mathsf{sk}\xleftarrow{\$}\{0,1\}^{256}$ is drawn
uniformly, so $\mathrm{Adv}^{\mathrm{PRF}}_{\mathrm{HMAC\text{-}SHA256}}(\mathcal{B}_1')$
is already the conventional uniform-key advantage the literature states,
and the term $q_F\cdot2^{-H_\infty(\mathbf{k})}$ is dropped entirely rather
than merely bounded. Second, the reduction $\mathcal{B}_1'$ constructed in
Lemma~\ref{lem:prf-hop}'s proof is instantiated exactly as
Lemma~\ref{lem:v8-cvf13} describes: $\mathcal{B}_1'$ samples its own
$\mathbf{k}'$ independently via $\mathsf{KeyGen}_{v8}$'s prime-sampling
step to run the arithmetic layer, and forwards every domain-derivation
call (including the synthetic-nonce domain \texttt{0x0A}) to its PRF
oracle --- a construction that Lemma~\ref{lem:v8-cvf13} shows is
distributed identically to the real/ideal games regardless of
$\mathcal{B}_1'$ not knowing $\mathsf{sk}^*$, closing the simulation gap
of Remark~\ref{rem:cvf13} that the analogous v7 reduction is subject to.
The nonce-collision term $q^2/2^{128}$ carries over unchanged: by
Lemma~\ref{lem:v8-cvf3}, two queries collide in synthetic nonce only with
probability at most
$\mathrm{Adv}^{\mathrm{PRF}}_{\mathrm{HMAC\text{-}SHA256}}(\mathcal{B}_1')+2^{-128}$,
which is already absorbed into the stated PRF term and the $q^2/2^{128}$
birthday bound over $q$ queries respectively --- no additional term is
introduced. The hiding argument (Lemma~\ref{lem:hiding}) and the
traffic-analysis argument (Lemma~\ref{lem:tar}) apply verbatim with
$\mathsf{kb}\to\mathsf{sk}$, since neither depends on how the nonce is
generated, only on it being fresh (which the collision bound above
supplies).
\end{proof}

\begin{corollary}[INT-CTXT and IND-CCA, v8]\label{cor:v8-security}
Let $\Pi_{v8}$ be as in Theorem~\ref{thm:ind-cpa-v8}. For any PPT adversary
making at most $q$ encryption queries, $q_v$ (resp.\ $q_d$)
verification/decryption queries, there exist PRF adversaries
$\mathcal{B}_2'$ (resp.\ $\mathcal{B}_1',\mathcal{B}_2'$) with similar
running time such that
\[
  \mathrm{Adv}^{\mathrm{INT\text{-}CTXT}}_{\Pi_{v8}}(\mathcal{A})
  \leq
  \mathrm{Adv}^{\mathrm{PRF}}_{\mathrm{HMAC\text{-}SHA256}}(\mathcal{B}_2')
  + \frac{q_v}{2^{256}},
  \qquad
  \mathrm{Adv}^{\mathrm{IND\text{-}CCA}}_{\Pi_{v8}}(\mathcal{A})
  \leq
  \mathrm{Adv}^{\mathrm{PRF}}_{\mathrm{HMAC\text{-}SHA256}}(\mathcal{B}_1')
  +
  \mathrm{Adv}^{\mathrm{PRF}}_{\mathrm{HMAC\text{-}SHA256}}(\mathcal{B}_2')
  + \frac{q^2}{2^{128}}
  + \frac{q_d}{2^{256}},
\]
both against the standard, uniform-key HMAC-SHA256 PRF assumption, with no
key-entropy term.
\end{corollary}

\begin{proof}
Theorem~\ref{thm:int-ctxt}'s proof (Case~1/Case~2) and
Theorem~\ref{thm:ind-cca}'s Bellare--Namprempre composition
\cite{bellare2000ae} depend on the shared wire-format/domain-separation
machinery of Section~\ref{sec:construction} (Remark~\ref{rem:domsep}'s
injectivity argument, the fixed-width token encoding, and the
domain-\texttt{0x03} tag) and on Theorem~\ref{thm:ind-cpa}'s IND-CPA bound
as a black box; neither proof inspects the key-generation or
nonce-derivation procedure beyond using it as an oracle. Substituting
$\mathsf{sk}$ for $\mathsf{kb}$ throughout (as in
Theorem~\ref{thm:ind-cpa-v8}'s proof) and invoking
Theorem~\ref{thm:ind-cpa-v8} in place of Theorem~\ref{thm:ind-cpa} at the
one point (Theorem~\ref{thm:ind-cca}'s ``Game $H_1$ reduces to IND-CPA''
step) where the IND-CPA bound is used, reproduces both proofs verbatim
with the key-guessing term dropped, exactly as in
Theorem~\ref{thm:ind-cpa-v8}.
\end{proof}

\begin{remark}[Scope and residual]
\label{rem:v8-scope}
Theorem~\ref{thm:ind-cpa-v8} and Corollary~\ref{cor:v8-security} ---
assembled from Lemmas~\ref{lem:v8-cvf3}--\ref{lem:v8-cvf13} --- close CVF3,
CVF8, and CVF13, and are the security statements this paper recommends new
deployments rely on. The v7 wire format, key schedule, and their
previously-documented residuals (Remarks~\ref{rem:cvf8-kdf},
\ref{rem:cvf8-kdf-erratum}, \ref{rem:cvf13}) are unchanged and remain
available, fully specified and separately proved
(Section~\ref{sec:algorithm-triple}, Theorem~\ref{thm:ind-cpa}), for
backward compatibility only; v7 callers who need the stronger guarantees
above must migrate to $\mathsf{KeyGen}_{v8}$/$\mathsf{Enc}_{v8}$/
$\mathsf{Dec}_{v8}$. \textbf{This migration is the paper's recommended
default as of this version.} The v8 schedule is implemented in
\texttt{rust/src/lib.rs} (\texttt{generate\_v8\_key},
\texttt{encrypt\_bytes\_v8}, \texttt{decrypt\_bytes\_v8}) with unit tests
covering round-trip correctness, tamper/AAD-mismatch rejection,
same-message determinism (the disclosed MRAE trade-off), and
distinct-message nonce distinctness. The Python (\texttt{napqes.py}) and C
(\texttt{C/napqes.c}) reference implementations have not yet been ported
to v8; this is tracked in \texttt{docs/CAVEATS.md} as a follow-up so all
three languages offer the misuse-resistant path.
\end{remark}

% ─────────────────────────────────────────────────────────────────────────────
\section{Security Analysis}
\label{sec:security}
% ─────────────────────────────────────────────────────────────────────────────

Theorem~\ref{thm:ind-cpa} (IND-CPA), the IND-CCA discussion of
Section~\ref{sec:ind-cca}, and the INT-CTXT theorem below are all stated
against the formal $(\mathsf{KeyGen}, \mathsf{Enc}, \mathsf{Dec})$ triple of
Definition~\ref{def:aead-triple} --- \textbf{the legacy v7 single-secret
construction} (Section~\ref{sec:algorithm-triple}) --- and assume the
correctness condition of Definition~\ref{def:correctness}; the
corresponding v8 statements, Theorem~\ref{thm:ind-cpa-v8} and
Corollary~\ref{cor:v8-security} (Section~\ref{sec:v8}), are proved by
direct reduction to the theorems in this section and carry the stronger,
key-entropy-free guarantee that is this paper's recommended default.
Together, INT-CTXT (ciphertext integrity
against an adversary with encryption-oracle access) and the underlying
HMAC-SHA256 tag's collision/forgery resistance are the EUF-CMA-style
unforgeability goal referenced informally in Section~\ref{sec:intro} —
NAPQES does not define a standalone MAC verification key separate from
$\mathsf{kb}$, so EUF-CMA unforgeability of the tag and INT-CTXT
unforgeability of the ciphertext coincide here. Arguments below that elide
$N$ or $A$ (e.g.\ writing $\mathsf{Enc}(\mathbf{k}, m)$) do so only where the
elided argument is fixed or immaterial to the hybrid being analysed;
the fully-typed calls are always as given in Definition~\ref{def:aead-triple}.

\subsection{IND-CPA Security}

\begin{remark}[Domain Separation of HMAC Inputs]\label{rem:domsep}
As of the CVF2 fix, every internal NAPQES derivation invokes
$F(\mathsf{kb},\cdot)$ with an input of the single, uniform shape
\[
  d \,\|\, N \,\|\, \mathsf{ctx},
\]
where $d \in \{\texttt{0x00},\ldots,\texttt{0x09}\}$ is a 1-byte domain
separator, $N$ is the 16-byte nonce (always at byte offset $1..16$), and
$\mathsf{ctx}$ is either a fixed-width positional counter (domains
$0,1,4,5$: 5\,B; domains $6,7$: 4\,B; domain $2$: empty) or a
length-prefixed, AAD-binding suffix (domains $3,8,9$:
$\mathrm{be4}(|A|) \| A \| \cdot$).

\noindent
\textbf{Cross-domain distinctness (unconditional).}
Two inputs $d_1\|N_1\|\mathsf{ctx}_1$ and $d_2\|N_2\|\mathsf{ctx}_2$ with
$d_1 \neq d_2$ differ at byte position~$0$ and are therefore always
distinct, regardless of $N$ or $\mathsf{ctx}$ — no probabilistic argument
is required.

\noindent
\textbf{Intra-domain injectivity.}
Fix a domain $d$. Bytes $1..16$ of every input to $\mathrm{Derive}_d$ are
the nonce $N$, so for a fixed $N$ the map $\mathsf{ctx} \mapsto d\|N\|\mathsf{ctx}$
is injective as long as $\mathsf{ctx} \mapsto \mathsf{ctx}$ itself is
injective. For domains $0,1,4,5,6,7$, $\mathsf{ctx}$ is a fixed-width
big-endian counter, so this holds trivially. For domains $3,8,9$,
$\mathsf{ctx} = \mathrm{be4}(|A|) \| A \| \mathsf{rest}$: because $\mathrm{be4}(|A|)$
fixes $|A|$, the split between $A$ and $\mathsf{rest}$ is unambiguous, so
distinct $(A, \mathsf{rest})$ pairs never collide — the length prefix makes
the AAD-binding encoding injective. Consequently, for a fixed nonce $N$,
distinct semantically-intended derivations under the same domain always
query $F(\mathsf{kb},\cdot)$ at distinct inputs, and (by the cross-domain
argument) derivations under different domains never collide either. No
cross-group probabilistic term (as required under the pre-CVF2 mixed
nonce-first/domain-first layout) is needed.

In the remainder of Section~\ref{sec:security} we treat all
$F(\mathsf{kb},\cdot)$ evaluations as queries to a single PRF; the
domain-separation argument above guarantees that semantically distinct
derivations always query the PRF at distinct inputs.
\end{remark}

\begin{remark}[Audit finding CVF8: the bound must depend on key entropy]
\label{rem:cvf8}
The original statement of Theorem~\ref{thm:ind-cpa} bounded
$\mathrm{Adv}^{\mathrm{IND\text{-}CPA}}_{\mathrm{NAPQES}}(\mathcal{A})$
purely in terms of $\mathrm{Adv}^{\mathrm{PRF}}_{\mathrm{HMAC\text{-}SHA256}}(\mathcal{B}_1)$
and a nonce-collision term, with no term depending on $K$ or $|\mathcal{P}|$.
This is unsound: $\mathsf{kb}$ is not a uniformly random bit-string, it is
$\mathsf{key\_bytes}(\mathbf{k})$ for $\mathbf{k}$ drawn from the
$2^{H_\infty(\mathbf{k})}$-point prime-tuple distribution of the Key
subsection, and the standard HMAC-SHA256 PRF conjecture is stated for a
uniformly random key. An adversary can always recover $\mathbf{k}$ (and
thereafter win the IND-CPA game outright) by exhaustively evaluating
$\mathrm{HMAC}(\cdot, \cdot)$ offline against a list of candidate keys, at
cost $\approx 2^{H_\infty(\mathbf{k})}$ evaluations in the worst case; this
attack is entirely outside the uniform-key PRF game and so is invisible to
a bound stated only in terms of $\mathrm{Adv}^{\mathrm{PRF}}_{\mathrm{HMAC\text{-}SHA256}}$.
For small $K$ (e.g.\ $K=1$, $H_\infty(\mathbf{k})\approx 19$--$20$ bits)
this search is trivial, so the original theorem was false as stated for
admissible small key sizes, and imprecise even at the paper's default
$K=10$. The restated theorem below adds the missing key-guessing term
$q_F \cdot 2^{-H_\infty(\mathbf{k})}$ (Lemma~\ref{lem:key-guess}) and
Remark~\ref{rem:min-K} gives the minimum $K$ the theorem requires to meet
a target security level.
\end{remark}

\begin{remark}[Audit finding CVF9: INDCPA was never formally defined, and the challenge constraint was ill-specified]
\label{rem:cvf9}
The original text stated Theorem~\ref{thm:ind-cpa} and its proof against
``the standard IND-CPA experiment'' without ever giving a formal IND-CPA
definition (adversary, oracles, challenge, advantage) — the experiment was
only sketched inline inside the proof of Game~$G_0$ below. That sketch
then imposed a challenge constraint, ``$\mathcal{A}$ submits a challenge
pair $(m_0,m_1)$ of \emph{equal padded length}'', which is non-standard
(the textbook constraint is $|m_0|=|m_1|$, equal \emph{plaintext} length,
with no padding-related precondition) and did not state which of three
genuinely different quantities --- codepoint count, on-wire byte length,
or the padded bucket $B$ (Plaintext Padding subsection) --- ``length''
referred to. Left unresolved, this made the theorem statement not
well-defined, and let the hiding lemma's proof (Lemma~\ref{lem:hiding})
later invoke ``equal padded lengths by the IND-CPA requirement'' as an
unproved assumption to close a step that would otherwise be circular.
Definition~\ref{def:ind-cpa} below supplies the missing formal definition,
using the standard constraint $|m_0|=|m_1|$ with the length metric stated
explicitly (codepoints, matching how $\mathcal{M}$ is typed in
Definition~\ref{def:aead-triple} --- the only metric NAPQES's algorithm
triple is defined against); no ``equal padded length'' precondition is
imposed on $\mathcal{A}$. Remark~\ref{rem:equal-padded-derived} then shows
that equal padded length is a \emph{consequence} of this standard
constraint here (not an additional restriction smuggled in to make the
proof close), because the padding bucket $B$ is a deterministic function
of codepoint count alone. Remark~\ref{rem:cvf9-byte-metric} states the
residual scope limitation this leaves: the guarantee is with respect to
$\mathcal{M}$'s codepoint metric, and does not, by itself, extend to
external byte-string encodings of unequal codepoint count but equal byte
length (audit finding CAV-003).
\end{remark}

\begin{definition}[IND-CPA security]\label{def:ind-cpa}
Let $\Pi=(\mathsf{KeyGen},\mathsf{Enc},\mathsf{Dec})$ be the NAPQES
algorithm triple of Definition~\ref{def:aead-triple}, and let
$\mathsf{Encrypt}(\mathbf{k},A,M)\triangleq\mathsf{Enc}(\mathbf{k},N,A,M)$
for $N\xleftarrow{\$}\mathcal{N}$ be its randomized public encryption
oracle (Section~\ref{sec:algorithm-triple}). For a PPT adversary
$\mathcal{A}=(\mathcal{A}_1,\mathcal{A}_2)$, define the experiment
$\mathbf{Exp}^{\mathrm{ind\text{-}cpa}}_{\Pi,\mathcal{A}}(\lambda)$:
\begin{enumerate}
  \item $\mathbf{k}\gets\mathsf{KeyGen}(1^\lambda)$.
  \item $\mathcal{A}_1$ is given adaptive oracle access to
        $\mathsf{Encrypt}(\mathbf{k},\cdot,\cdot)$ and, after at most $q$
        total queries (counted jointly with step 4), outputs
        $(A^*,m_0,m_1,\mathsf{st})$ with $m_0,m_1\in\mathcal{M}$ subject to
        exactly one constraint:
        \[
          |m_0| = |m_1|,
        \]
        where $|\cdot|$ denotes the number of Unicode codepoints of a
        message in $\mathcal{M}$ --- the metric $\mathcal{M}$ is typed
        against in Definition~\ref{def:aead-triple}. No other restriction
        (padded length, byte length, or otherwise) is imposed on
        $\mathcal{A}_1$'s choice of $m_0,m_1$.
  \item The challenger samples $b\xleftarrow{\$}\{0,1\}$ and returns
        $c^*\gets\mathsf{Encrypt}(\mathbf{k},A^*,m_b)$.
  \item $\mathcal{A}_2(\mathsf{st},c^*)$, with continued oracle access to
        $\mathsf{Encrypt}(\mathbf{k},\cdot,\cdot)$, outputs a bit $b'$.
  \item The experiment outputs $1$ (``$\mathcal{A}$ wins'') iff $b'=b$.
\end{enumerate}
The advantage of $\mathcal{A}$ is
\[
  \mathrm{Adv}^{\mathrm{IND\text{-}CPA}}_{\mathrm{NAPQES}}(\mathcal{A})
  \;=\;
  \Bigl|\Pr\bigl[\mathbf{Exp}^{\mathrm{ind\text{-}cpa}}_{\Pi,\mathcal{A}}(\lambda)=1\bigr]-\tfrac12\Bigr|,
\]
and $\Pi$ is IND-CPA secure if this quantity is negligible in $\lambda$ for
every PPT $\mathcal{A}$. This is the standard left-or-right IND-CPA notion
for a symmetric encryption scheme, stated against NAPQES's typed algorithm
triple with no bespoke preconditions; Theorem~\ref{thm:ind-cpa} is proved
against exactly this definition, and $q$ there is the total query count of
this definition.
\end{definition}

\begin{remark}[Equal padded length is derived, not assumed]
\label{rem:equal-padded-derived}
Definition~\ref{def:ind-cpa} imposes only the standard constraint
$|m_0|=|m_1|$ (equal codepoint count); it does not restrict $\mathcal{A}$
to equal \emph{padded} length. The Plaintext Padding subsection shows
$B=\max(16,2^{\lceil\log_2(n+1)\rceil})$ is a deterministic function of
$n=|P|$ (the codepoint count) alone, with no dependence on codepoint
\emph{values}. Consequently $|m_0|=|m_1|$ already forces $m_0$ and $m_1$
to share the same padding bucket $B$ and the same padded length $B+2$,
with certainty --- this is a logical consequence of
Definition~\ref{def:ind-cpa}'s standard constraint together with the
padding formula, not an independent assumption. Every subsequent
occurrence of ``$m_0,m_1$ have equal padded length'' below (Game $G_0$'s
challenge step, Lemma~\ref{lem:hiding}'s proof) refers to this derived
fact, closing the circularity audit finding CVF9 identified in the
original text, which instead stated equal padded length as an unproved
precondition and later invoked it as if it were the definition's own
requirement.
\end{remark}

\begin{remark}[Scope: the metric is codepoints, not on-wire bytes (residual of CVF9)]
\label{rem:cvf9-byte-metric}
Definition~\ref{def:ind-cpa} measures $|m_0|=|m_1|$ in Unicode codepoints
because that is the only metric $\mathcal{M}$ is formally typed against
(Definition~\ref{def:aead-triple}); this is not a free choice among
equally natural alternatives, since NAPQES's $\mathsf{Enc}$ takes a
codepoint sequence, not a byte string, as its message input. It follows
that Theorem~\ref{thm:ind-cpa} says nothing about a caller-level scenario
in which two externally byte-equal-length plaintexts (e.g.\ two UTF-8
byte strings of the same length) decode to \emph{different} codepoint
counts: such a pair falls outside $\mathcal{A}$'s admissible challenge set
entirely (they are not even a valid $(m_0,m_1)\in\mathcal{M}\times\mathcal{M}$
challenge unless re-expressed at the codepoint level), and if instead
compared as codepoint-level messages of unequal length, the padded bucket
$B$ --- and hence ciphertext length --- may differ between them,
consistent with the length-bucket leak already documented as CAV-003. A
caller relying on NAPQES to hide the distinction between two plaintexts
that are equal-length only in bytes, not in codepoints, is not covered by
Theorem~\ref{thm:ind-cpa} and should consult CAV-003 (Section~\ref{sec:caveats}).
\end{remark}

\begin{theorem}[IND-CPA]\label{thm:ind-cpa}
  Under the PRF assumption on HMAC-SHA256 (Remark~\ref{rem:prf-d} makes
  precise which key distribution this assumption is invoked against once
  the guessing term below has been paid for, and Remark~\ref{rem:cvf13}
  makes precise an additional simulation-gap caveat the reduction $\mathcal{B}_1$
  is subject to), the NAPQES block encryption
  scheme is IND-CPA secure (Definition~\ref{def:ind-cpa}) provided $K$ and
  $|\mathcal{P}|$ are large enough
  that $H_\infty(\mathbf{k})$ meets the target security level
  (Remark~\ref{rem:min-K}). Specifically, for any PPT adversary
  $\mathcal{A}$ making at most $q$ encryption queries and at most $q_F$
  offline evaluations of $\mathrm{HMAC}(\cdot,\cdot)$ under
  adversarially-chosen keys, there exists a PRF adversary $\mathcal{B}_1$
  with similar running time such that:
  \[
    \mathrm{Adv}^{\mathrm{IND\text{-}CPA}}_{\mathrm{NAPQES}}(\mathcal{A})
    \leq
    \mathrm{Adv}^{\mathrm{PRF}}_{\mathrm{HMAC\text{-}SHA256}}(\mathcal{B}_1)
    + \frac{q^2}{2^{128}}
    + q_F \cdot 2^{-H_\infty(\mathbf{k})}.
  \]
\end{theorem}

\begin{proof}
We prove the bound via a sequence of two hybrid games.

\paragraph{Game $G_0$ (Real Game).}
This is the standard IND-CPA experiment.  The challenger samples
$\mathbf{k}=(k_1,\ldots,k_K)$ uniformly from $K$-element ordered tuples
of distinct primes in~$\mathcal{P}$, computes $\mathsf{kb}=\mathsf{key\_bytes}(\mathbf{k})$,
and answers up to $q$ chosen-plaintext queries honestly using the real
$\operatorname{NAPQES}.\mathsf{Enc}$ algorithm (which calls
$F(\mathsf{kb},\cdot)=\operatorname{HMAC\text{-}SHA256}(\mathsf{kb},\cdot)$
for all internal derivations).  After the query phase,
$\mathcal{A}$ submits a challenge pair $(m_0,m_1)$ with $|m_0|=|m_1|$
(Definition~\ref{def:ind-cpa}) --- which, by
Remark~\ref{rem:equal-padded-derived}, already forces $m_0,m_1$ to have
equal padded length; no separate precondition is imposed here.
The challenger samples $b\xleftarrow{\$}\{0,1\}$, returns
$c^*\gets\operatorname{NAPQES}.\mathsf{Enc}(\mathbf{k},m_b)$, and
$\mathcal{A}$ wins if its output $b'=b$.

\paragraph{Game $G_1$ (PRF Replacement).}
Identical to $G_0$, except the function $F(\mathsf{kb},\cdot)$ is
replaced by a truly random function $\rho:\{0,1\}^*\to\{0,1\}^{256}$,
sampled uniformly at the start of the experiment and evaluated lazily.
Every call that would invoke $\operatorname{HMAC}(\mathsf{kb},x)$ now
returns $\rho(x)$ instead.

\begin{lemma}[Key-guessing bound (audit finding CVF8)]\label{lem:key-guess}
Let $q_F$ upper-bound the number of evaluations of
$\mathrm{HMAC}(\kappa,\cdot)$ that $\mathcal{A}$ performs \emph{offline},
under keys $\kappa$ of its own choosing, in addition to its $q$ online
encryption queries (this is exactly the resource the original theorem
statement left unbounded --- see Remark~\ref{rem:cvf8}). Let $\mathrm{Guess}$
denote the event that some such $\kappa$ equals the challenger's actual
$\mathsf{kb}=\mathsf{key\_bytes}(\mathbf{k})$. Since $\mathbf{k}$ is sampled
uniformly from the $2^{H_\infty(\mathbf{k})}$ ordered $K$-tuples of distinct
primes in $\mathcal{P}$ (Key subsection) independently of $\mathcal{A}$'s
(unbounded-time, but $q_F$-query) offline guessing strategy, a union bound
over the $q_F$ guesses gives
\[
  \Pr[\mathrm{Guess}] \;\leq\; q_F \cdot 2^{-H_\infty(\mathbf{k})}.
\]
If $\mathrm{Guess}$ occurs, $\mathcal{A}$ can compute
$F(\mathsf{kb},\cdot)$ directly and simulate the challenger perfectly,
winning $G_0$ with probability $1$; this is the exhaustive-key-search
attack the original, entropy-free bound could not see. We therefore only
bound $\mathcal{A}$'s advantage in $G_0$ conditioned on
$\overline{\mathrm{Guess}}$ via the PRF hop below, and add
$\Pr[\mathrm{Guess}]$ back in when combining the bounds.
\end{lemma}

\begin{lemma}\label{lem:prf-hop}
There exists a PRF adversary $\mathcal{B}_1$, running in time
$\mathcal{O}(T_{\mathcal{A}})$ with at most $q\cdot L_{\max}$ oracle calls
(where $L_{\max}$ is an upper bound on the number of HMAC evaluations per
encryption), such that
\[
  \bigl|\Pr[\mathcal{A}\text{ wins }G_0]-\Pr[\mathcal{A}\text{ wins }G_1]\bigr|
  \;\leq\;
  \mathrm{Adv}^{\mathrm{PRF}}_{F}(\mathcal{B}_1).
\]
\end{lemma}

\begin{proof}
$\mathcal{B}_1$ is given a function oracle $\mathcal{O}(\cdot)$ that is
either $F(\mathsf{kb},\cdot)$ (real) or $\rho(\cdot)$ (random).
$\mathcal{B}_1$ simulates the IND-CPA experiment for $\mathcal{A}$, answering
every internal HMAC call by forwarding to $\mathcal{O}$, and
(see Remark~\ref{rem:cvf13} for why this description alone is
incomplete, and how the remaining gap is resolved) samples $\mathbf{k}$
itself in order to run the arithmetic layer of $\mathrm{NAPQES}.\mathsf{Enc}$.
If $\mathcal{O}=F(\mathsf{kb},\cdot)$, the simulation is identical to $G_0$;
if $\mathcal{O}=\rho$, it is identical to $G_1$.
Hence $\mathcal{B}_1$'s distinguishing advantage equals
$|\Pr[\mathcal{A}\text{ wins }G_0]-\Pr[\mathcal{A}\text{ wins }G_1]|$,
conditioned throughout on $\overline{\mathrm{Guess}}$ (Lemma~\ref{lem:key-guess}).
\end{proof}

\begin{remark}[What $\mathrm{Adv}^{\mathrm{PRF}}_{\mathrm{HMAC\text{-}SHA256}}(\mathcal{B}_1)$ means here]
\label{rem:prf-d}
Conditioned on $\overline{\mathrm{Guess}}$, Lemma~\ref{lem:prf-hop} shows
that distinguishing $G_0$ from $G_1$ is exactly as hard as distinguishing
$F(\mathsf{kb},\cdot)$ from a random function \emph{for $\mathsf{kb}$ drawn
from the prime-tuple distribution $D$ of the Key subsection, conditioned on
not having been one of $\mathcal{A}$'s $q_F$ offline guesses} --- not for
$\mathsf{kb}$ drawn uniformly from $\{0,1\}^{40K}$. These are different
hypotheses: the HMAC-SHA256 PRF conjecture, as conventionally stated and
studied, concerns a uniformly random key. $D$ is a highly structured,
measure-zero subset of $\{0,1\}^{40K}$ (5-byte-aligned fields, each in a
prime-valued range, no repeats, order-independent as a set) with min-entropy
$H_\infty(\mathbf{k}) = \log_2|\mathcal{P}|!/(|\mathcal{P}|-K)! \ll 40K$.
We write $\mathrm{Adv}^{\mathrm{PRF}}_{\mathrm{HMAC\text{-}SHA256}}(\mathcal{B}_1)$
in Theorem~\ref{thm:ind-cpa} to denote the PRF advantage against keys drawn
from $D$ (conditioned on $\overline{\mathrm{Guess}}$) specifically --- a
strictly weaker, and far less studied, assumption than the uniform-key one
--- and we do not claim the standard uniform-key HMAC-SHA256 PRF assumption
directly implies this quantity is negligible. Section~\ref{sec:cvf8-fix}
discusses two ways to remove this residual non-standard assumption and
restore a reduction to the standard, uniform-key conjecture.
\end{remark}

\begin{remark}[Audit finding CVF13: $\mathcal{B}_1$/$\mathcal{B}_2$ cannot simulate NAPQES encryption from oracle access alone]
\label{rem:cvf13}
Lemma~\ref{lem:prf-hop}'s proof (and, symmetrically, $\mathcal{B}_2$'s
construction in the proof of Theorem~\ref{thm:int-ctxt}, Case~1) states
that the reduction ``simulates the \dots\ experiment for $\mathcal{A}$,
answering every internal HMAC call by forwarding to $\mathcal{O}$/$f$.''
This is necessary but not sufficient: to answer $\mathcal{A}$'s
encryption queries at all, the reduction must run
$\mathrm{NAPQES}.\mathsf{Enc}(\mathbf{k},\cdot)$, and that algorithm uses
$\mathbf{k}$ itself \emph{outside} of any HMAC call --- the token emission
step computes $c\cdot k_j + a$ (Section~\ref{sec:construction}) and the
admissible addend range $[1,k_j-1]$
(Definitions~\ref{def:noiseaddend}--\ref{def:realaddend}) both depend on
the specific prime $k_j$, not merely on $\mathrm{HMAC}(\mathsf{kb},\cdot)$'s
output. Forwarding HMAC calls to an oracle $\mathcal{O}$ does not supply
$\mathbf{k}$: in the standard PRF game $\mathcal{O}$'s key is sampled
uniformly and kept secret from the distinguisher, so $\mathcal{B}_1$
(resp.\ $\mathcal{B}_2$) has no way to extract $\mathbf{k}$ from oracle
access, and cannot run the arithmetic layer as written.

The reduction cannot be repaired by having $\mathcal{B}_1$ sample its own
$\mathbf{k}'$ locally for the arithmetic layer: whenever
$\mathcal{O}=F(\mathsf{kb},\cdot)$ (the real-world branch), there is no
reason $\mathsf{kb}$ equals $\mathsf{key\_bytes}(\mathbf{k}')$ for the
\emph{specific} $\mathbf{k}'$ that $\mathcal{B}_1$ chose --- $\mathbf{k}'$
and the PRF challenger's hidden key are independent random variables. The
resulting simulation would answer $\mathcal{A}$'s queries with tokens
computed under $\mathbf{k}'$ but tags/addends/positions derived from
$F(\mathsf{kb},\cdot)$ for an unrelated $\mathbf{k}$: an internally
inconsistent hybrid that is a different scheme from
$\mathrm{NAPQES}.\mathsf{Enc}(\mathbf{k}',\cdot)$ for any single key, and
so does not establish $\Pr[\mathcal{A}\text{ wins }G_0]$ as claimed.
Symmetrically, if $\mathcal{B}_1$ instead supplies its own $\mathbf{k}'$
\emph{to} the PRF challenger (a ``chosen-key'' PRF game, so that
$\mathcal{O}$'s hidden key really is $\mathsf{key\_bytes}(\mathbf{k}')$),
$\mathcal{B}_1$ already knows the tested key and can compute
$F(\mathsf{key\_bytes}(\mathbf{k}'),x)$ itself, making its distinguishing
advantage against $\mathcal{O}$ trivially $\approx 1$ regardless of
whether $\mathcal{O}$ is real or random --- this shows nothing, since the
premise of a PRF game (the key is hidden from the distinguisher) is
violated by construction whenever the same key material must also be
disclosed, in the clear, to run the arithmetic layer.

This defect is shared verbatim by the INT-CTXT reduction $\mathcal{B}_2$
(Theorem~\ref{thm:int-ctxt}'s proof, Case~1), which likewise must run
$\mathrm{NAPQES}.\mathsf{Enc}$ to answer $\mathcal{A}$'s $q$ encryption
queries before evaluating any forgery, and it therefore propagates into
the IND-CCA bound (Theorem~\ref{thm:ind-cca}), which is assembled from
$\mathcal{B}_1$'s and $\mathcal{B}_2$'s advantages via the
Bellare--Namprempre composition~\cite{bellare2000ae}: neither advantage
is actually established for the real NAPQES scheme as the reductions are
currently written, so the composed IND-CCA bound rests on the same
unproven premise. This is a distinct defect from CVF8/Remark~\ref{rem:prf-d}
(which concerns the \emph{key distribution} the PRF assumption is stated
over) --- CVF13 concerns whether the reduction can be run \emph{at all},
independent of which key distribution is assumed.

\textbf{Resolution.} Remark~\ref{rem:cvf8-kdf}'s KDF-subkey design ---
deriving an independent, uniformly-distributed
$\mathsf{sk}=\mathrm{HKDF}(\mathsf{kb})$ and keying every domain
derivation with $\mathsf{sk}$ in place of $\mathsf{kb}$, while
$\mathbf{k}$ itself is used only for the (public, non-HMAC) arithmetic
layer --- resolves both defects simultaneously: once $\mathsf{sk}$ is
architecturally independent of $\mathbf{k}$, $\mathcal{B}_1$/$\mathcal{B}_2$
may legitimately sample $\mathbf{k}$ locally (it is no longer correlated
with the tested key at all, so no ``chosen-key'' game or coincidence is
needed) while forwarding every domain-derivation call to an external
oracle keyed by the PRF challenger's independent, hidden $\mathsf{sk}$.
Absent that change, the reductions as written require an explicit,
additional \emph{key-indistinguishability} step that this proof does not
supply: a demonstration that a simulator holding $\mathbf{k}$ in the
clear (for arithmetic) still cannot distinguish
$\mathrm{HMAC}(\mathsf{key\_bytes}(\mathbf{k}),\cdot)$ from random using
that knowledge alone --- i.e., that $\mathbf{k}$'s use in the arithmetic
layer leaks nothing that helps predict HMAC outputs under
$\mathsf{key\_bytes}(\mathbf{k})$. We do not supply that step here: it is
not implied by the standard HMAC-PRF conjecture (which is silent about a
distinguisher that possesses the key), and we are not aware of a
suitable weaker primitive-level assumption. Theorem~\ref{thm:ind-cpa},
Theorem~\ref{thm:int-ctxt}, and Theorem~\ref{thm:ind-cca} should
therefore be read as conditioned on this gap until either the
KDF-subkey design lands or the missing key-indistinguishability step is
supplied; we flag this as an open residual, cross-referenced from
Section~\ref{sec:cvf8-fix} and logged in Section~\ref{sec:caveats}.
\end{remark}

\paragraph{Nonce collision probability.}
Let $\mathrm{Coll}$ denote the event that at least two of the $q+1$ nonces
generated across the $q$ CPA oracle calls and the challenge encryption are
equal.  Each nonce is drawn independently and uniformly from $\{0,1\}^{128}$,
so by the union bound:
\[
  \Pr[\mathrm{Coll}]
  \;\leq\;
  \binom{q+1}{2}\cdot 2^{-128}
  \;\leq\;
  \frac{q^2}{2^{128}}.
\]

\begin{lemma}\label{lem:hiding}
In $G_1$ conditioned on $\overline{\mathrm{Coll}}$:
\[
  \Pr\bigl[\mathcal{A}\text{ wins }G_1\;\big|\;\overline{\mathrm{Coll}}\bigr]
  \;=\;\tfrac{1}{2}.
\]
\end{lemma}

\begin{proof}
Fix any adversary view $V$ arising from the CPA query phase, which uses
nonces $N_1,\ldots,N_q$.  Let $N^*$ be the challenge nonce.
Under $\overline{\mathrm{Coll}}$, $N^*\notin\{N_1,\ldots,N_q\}$.

In $G_1$, every challenge-phase derivation uses an input prefixed by the
domain byte $d$ and then $N^*$
(e.g., $\rho(\texttt{0x00}\|N^*\|\mathrm{be5}(j))$,
$\rho(\texttt{0x01}\|N^*\|\mathrm{be5}(j))$, $\rho(\texttt{0x02}\|N^*)$,
and so on for domains \texttt{0x03}--\texttt{0x07}).
All CPA-phase evaluations used inputs containing some $N_i\neq N^*$ at the
same fixed byte $1..16$ offset.
Since $\rho$ is a truly random function, its outputs on distinct inputs are
independent, uniform $256$-bit values.
Under $\overline{\mathrm{Coll}}$, the challenge inputs are disjoint from
all CPA-phase inputs, so all challenge $\rho$-outputs are fresh, uniform,
and independent of~$V$.

In particular, the keystream masking the token blob (domain~\texttt{0x07}) is:
\[
  \mathrm{ks}(N^*)
  = \rho(\texttt{0x07}\|N^*\|\mathrm{be4}(0))
    \;\|\;
    \rho(\texttt{0x07}\|N^*\|\mathrm{be4}(1))
    \;\|\;\cdots
  \quad(\text{truncated to }|B|\text{ bytes}),
\]
where $B$ denotes the fixed-width-encoded token sequence (Section~\ref{sec:wire-format-v7}).

\smallskip
\noindent\textbf{Independence of $\mathrm{ks}(N^*)$ from $B$ (OTP argument).}
$B$ is assembled entirely from $\rho$-outputs queried under domains
$\{\texttt{0x00},\texttt{0x01},\texttt{0x02},\texttt{0x04},\texttt{0x05},\texttt{0x06}\}$:
\begin{itemize}
  \item Noise/real positions: $\rho(\texttt{0x00}\|N^*\|\mathrm{be5}(\cdot))$
  \item Real/noise addends: $\rho(\texttt{0x01}\|N^*\|\mathrm{be5}(\cdot))$,
        $\rho(\texttt{0x05}\|N^*\|\mathrm{be5}(\cdot))$
  \item Noise probability: $\rho(\texttt{0x02}\|N^*)$
  \item Noise codepoints: $\rho(\texttt{0x04}\|N^*\|\mathrm{be5}(\cdot))$
  \item Padding: $\rho(\texttt{0x06}\|N^*\|\mathrm{be4}(\cdot))$
\end{itemize}
Each keystream block uses input $\texttt{0x07}\|N^*\|\mathrm{be4}(j)$, domain
byte \texttt{0x07}.
By Remark~\ref{rem:domsep} (cross-domain distinctness, unconditional — every
input above shares the single uniform $d\|N\|\mathsf{ctx}$ shape), no
keystream input $\texttt{0x07}\|N^*\|\mathrm{be4}(j)$ equals any $B$-construction
input: they differ at byte position~$0$ (\texttt{0x07} vs.\ \texttt{0x00}--\texttt{0x06}).
No probabilistic cross-group term is needed — the argument is exact, not an
estimate over $q$ queries (this closes audit finding CVF2's cross-reference,
issue \#848, which tracked the removal of that probabilistic term).
Since $\rho$ is a truly random function, outputs at \emph{distinct} inputs
are mutually independent; therefore
$\mathrm{ks}(N^*)$ is statistically independent of~$B$.

Consequently, the one-time-pad identity applies: for any fixed deterministic $B$,
\[
  \tilde{B}
  = B \oplus \mathrm{ks}(N^*)
  \;\overset{\mathrm{d}}{=}\;
  \mathrm{Uniform}\bigl(\{0,1\}^{8|B|}\bigr).
\]

\smallskip
\noindent\textbf{Equal byte-length of $B$ across $m_0, m_1$ (CVF1 fix).}
Write $B = \mathrm{be8}(\mathrm{token}[0]) \| \cdots \| \mathrm{be8}(\mathrm{token}[N-1])$, so
$|B| = 8N$ (Section~\ref{sec:wire-format-v7}). The token count $N$ (real $+$
noise tokens) is determined entirely by: (i) the padded codepoint count,
which is equal for $m_0$ and $m_1$ because $|m_0|=|m_1|$
(Definition~\ref{def:ind-cpa}) forces equal padded length
(Remark~\ref{rem:equal-padded-derived} --- a derived fact, not an assumed
requirement), and (ii) the noise-position oracle
$\rho(N^*\|\texttt{0x00}\|\mathrm{be5}(\cdot))$, which depends only on
$N^*$ and $\mathsf{kb}$ (via $\rho$) --- \emph{not} on the codepoint values
being encrypted. Since the challenger uses the same $N^*$ and the same key
regardless of which of $m_0, m_1$ is encrypted, the sequence of
noise/real decisions at each $ct\_pos$ is identical in both cases, so $N$ ---
and hence $|B| = 8N$ --- is identical for $m_0$ and $m_1$. Therefore $\tilde{B}$
is uniformly distributed over strings of a length that is a deterministic
function of the (equal) padded length alone, independently of the
plaintext bit~$b$ and the view~$V$.

\begin{remark}[Erratum: the v6 argument was false]
\label{sec:hiding-lemma-erratum}
In the retired v6 construction, $B$ was instead
$\mathrm{varint}(\mathrm{token}[0]) \| \cdots \| \mathrm{varint}(\mathrm{token}[N-1])$,
and $\mathrm{varint}(\cdot)$ is \emph{not} a fixed-width encoding: its
byte-length grows with $\mathrm{token} = c \cdot k + a$, hence with the
plaintext codepoint $c$. Equal padded length only fixes $N$; it does not fix
$\sum_i \lvert \mathrm{varint}(\mathrm{token}[i]) \rvert$, which depends on
codepoint \emph{values}. The original v6 proof asserted ``$B$ is the same
length for both plaintexts $m_0, m_1$ (they have equal padded lengths by the
IND-CPA requirement)'' without justification --- this step is false, as
demonstrated by audit finding CVF1 with an explicit distinguisher
($m_0 = $ U+0001 repeated vs.\ $m_1 = $ U+FFFF repeated, equal padded
length, ciphertexts of different length with overwhelming probability, IND-CPA
advantage $\approx 1$). The fixed-width encoding of
Section~\ref{sec:wire-format-v7} removes the dependency of $|B|$ on
codepoint values, which is what restores the step above.
\end{remark}

\smallskip
\noindent\textbf{Freshness of the authentication tag.}
The tag is computed as
$T = \rho(\texttt{0x03}\|N^*\|\mathrm{be4}(|aad|)\|aad\|\tilde{B})$
(CVF2 fix: domain \texttt{0x03} uses the same unified $d\|N\|\mathsf{ctx}$
layout as every other domain; see Table~\ref{tab:domains}).
The tag input has first byte $\texttt{0x03}$ and, at bytes $1..16$, $N^*$.
All CPA-phase tag inputs are of the form
$\texttt{0x03}\|N_i\|\mathrm{be4}(|aad_i|)\|aad_i\|\tilde{B}_i$ with $N_i \neq N^*$;
since $N^*$ occupies the same fixed byte offset in both, these strings differ
unconditionally whenever $N_i \neq N^*$ --- which holds under
$\overline{\mathrm{Coll}}$, with no further probabilistic argument needed.
By Remark~\ref{rem:domsep} (cross-domain distinctness), the tag input also
differs unconditionally from every other domain's $\rho$-input (they differ
at byte position~$0$).
Therefore $T$ is a fresh, uniformly distributed $32$-byte value independent
of $b$ and~$V$.

Consequently, $c^*=(N^*,\tilde{B},T)$ is jointly uniformly distributed
over its domain, independently of~$b$.  The adversary gains no information
about $b$ from $c^*$, so
$\Pr[\mathcal{A}\text{ wins }G_1\mid\overline{\mathrm{Coll}}]=\tfrac{1}{2}$.
\end{proof}

\begin{remark}[Scope of Lemma~\ref{lem:hiding} (audit finding CVF4)]
\label{rem:hiding-scope}
The proof above never inspects the internal structure of $B$: the
OTP argument only uses that $\mathrm{ks}(N^*)$ (domain \texttt{0x07}) is
independent of $B$, for \emph{any} fixed deterministic $B$. The prime
multiplication $c \cdot k + a$ and the noise tokens assembled into $B$
(domains \texttt{0x00}, \texttt{0x01}, \texttt{0x02}, \texttt{0x04},
\texttt{0x05}, \texttt{0x06}) therefore play no role in this
confidentiality argument: per-message content confidentiality (IND-CPA)
reduces entirely to domain \texttt{0x07} being a secure keyed PRF and to
the domain-\texttt{0x03} tag for integrity. This does not make the
prime/noise layer redundant --- it establishes a different, and
independent, property; see Lemma~\ref{lem:tar}.
\end{remark}

\subsection{Traffic-Analysis Resistance (a property independent of Lemma~\ref{lem:hiding})}
\label{sec:traffic-analysis}

\begin{lemma}[Ciphertext length is decorrelated from plaintext content]
\label{lem:tar}
For any two plaintexts $m_0, m_1$ of equal padded codepoint length $n$,
encrypted under the same key $\mathsf{kb}$ and the same nonce $N$, the
ciphertext byte-length $|C| = 16 + 8N_{\mathrm{tok}} + 32$ (where
$N_{\mathrm{tok}}$ is the real-plus-noise token count) is identical for
$m_0$ and $m_1$, and depends only on $(\mathsf{kb}, N, n)$ --- never on the
codepoint values of $m_0$ or $m_1$ --- even in a hypothetical world where
domain \texttt{0x07} is fully distinguishable from random (i.e.\ this
holds independently of Lemma~\ref{lem:hiding}).
\end{lemma}

\begin{proof}
The token-construction loop (Section~\ref{sec:wire-format-v7}) advances
$ct\_pos$ by querying $\mathrm{is\_noise}(ct\_pos) = \rho(\texttt{0x00}\|N\|\mathrm{be5}(ct\_pos))$
until $n$ real-token slots have been consumed; this sequence of oracle
calls reads only $(\mathsf{kb}, N, ct\_pos)$ and never the codepoint value
being emitted at a real-token slot. Consequently $N_{\mathrm{tok}}$ is a
deterministic function of $(\mathsf{kb}, N, n)$ alone, identical for $m_0$
and $m_1$ since both have padded length $n$ under the same $(\mathsf{kb},
N)$. The fixed-width encoding $|B| = 8 N_{\mathrm{tok}}$
(Section~\ref{sec:wire-format-v7}, CVF1 fix) fixes the blob length
regardless of token \emph{values}, and domain \texttt{0x07} XOR-masks $B$
without changing its length ($|\widetilde{B}|=|B|$). The tag $T$ is a
fixed 32-byte HMAC output (Table~\ref{tab:domains}). Hence
$|C| = 16 + 8N_{\mathrm{tok}} + 32$ is identical for $m_0, m_1$ and is a
function of $(\mathsf{kb}, N, n)$ only, regardless of whether domain
\texttt{0x07} is secure.
\end{proof}

\begin{remark}[Scope: this is not a confidentiality claim]
Lemma~\ref{lem:tar} says nothing about whether ciphertext \emph{content}
reveals $m_0$ vs.\ $m_1$ --- that is Lemma~\ref{lem:hiding} and Remark~\ref{rem:hiding-scope}
above, which depend entirely on domain \texttt{0x07} and are indifferent to
$B$'s internal structure. Lemma~\ref{lem:tar} instead formalises the
noise/prime layer's actual, distinct contribution (audit finding CVF4):
ciphertext length does not correlate with plaintext content, a property
AES-GCM and ChaCha20-Poly1305 ciphertexts do not have, since their length
is an exact linear function of plaintext length. The $2^{197}$--$2^{257}$
prime-tuple key-space figures (Table~\ref{tab:comparison}) are a
key-recovery-resistance statement about the keyed-HMAC construction as a
whole, not a statement about what the token structure contributes to
either Lemma~\ref{lem:hiding} or Lemma~\ref{lem:tar} individually.
\end{remark}

\paragraph{Combining the bounds.}
By Lemma~\ref{lem:key-guess}, $\Pr[\mathrm{Guess}]\leq q_F\cdot 2^{-H_\infty(\mathbf{k})}$,
and $\mathcal{A}$ wins $G_0$ with probability at most $1$ whenever
$\mathrm{Guess}$ occurs, so
\[
  \Pr[\mathcal{A}\text{ wins }G_0]
  \;\leq\;
  q_F\cdot 2^{-H_\infty(\mathbf{k})}
  \;+\;
  \Pr[\mathcal{A}\text{ wins }G_0 \mid \overline{\mathrm{Guess}}].
\]
Conditioned throughout on $\overline{\mathrm{Guess}}$, the $G_0$/$G_1$ hop
(Lemma~\ref{lem:prf-hop}, Remark~\ref{rem:prf-d}) and the hiding lemma
proceed exactly as in the original argument:
\begin{align*}
  \Pr[\mathcal{A}\text{ wins }G_1 \mid \overline{\mathrm{Guess}}]
  &= \Pr[\mathrm{Coll}]\cdot\Pr[\text{wins}\mid\mathrm{Coll},\overline{\mathrm{Guess}}]
     +\Pr[\overline{\mathrm{Coll}}]\cdot\Pr[\text{wins}\mid\overline{\mathrm{Coll}},\overline{\mathrm{Guess}}]\\
  &\leq \Pr[\mathrm{Coll}]+\tfrac{1}{2}
   \;\leq\; \frac{q^2}{2^{128}}+\tfrac{1}{2},
\end{align*}
so that
\[
  \Bigl|\Pr[\mathcal{A}\text{ wins }G_0 \mid \overline{\mathrm{Guess}}]-\tfrac{1}{2}\Bigr|
  \;\leq\;
  \bigl|\Pr[G_0\mid\overline{\mathrm{Guess}}]-\Pr[G_1\mid\overline{\mathrm{Guess}}]\bigr|
  +\Bigl|\Pr[\mathcal{A}\text{ wins }G_1\mid\overline{\mathrm{Guess}}]-\tfrac{1}{2}\Bigr|
  \;\leq\;
  \mathrm{Adv}^{\mathrm{PRF}}_{F}(\mathcal{B}_1) + \frac{q^2}{2^{128}}.
\]
Therefore:
\begin{align*}
  \mathrm{Adv}^{\mathrm{IND\text{-}CPA}}_{\mathrm{NAPQES}}(\mathcal{A})
  &= \bigl|\Pr[\mathcal{A}\text{ wins }G_0]-\tfrac{1}{2}\bigr| \\
  &\leq q_F\cdot 2^{-H_\infty(\mathbf{k})}
       + \Bigl|\Pr[\mathcal{A}\text{ wins }G_0 \mid \overline{\mathrm{Guess}}]-\tfrac{1}{2}\Bigr| \\
  &\leq \mathrm{Adv}^{\mathrm{PRF}}_{F}(\mathcal{B}_1) + \frac{q^2}{2^{128}}
       + q_F\cdot 2^{-H_\infty(\mathbf{k})}.
    \qedhere
\end{align*}
\end{proof}

\subsection{CVF8: Minimum Key Size and Removing the Residual Non-Standard Assumption}
\label{sec:cvf8-fix}

\begin{remark}[Minimum $K$ for a target security level]
\label{rem:min-K}
Since $|\mathcal{P}|\gg K$, $H_\infty(\mathbf{k}) = \log_2|\mathcal{P}|!/(|\mathcal{P}|-K)!
\approx K\log_2|\mathcal{P}| \approx 19.16\,K$ bits for $|\mathcal{P}|\approx
586{,}000$. Theorem~\ref{thm:ind-cpa}'s key-guessing term
$q_F\cdot 2^{-H_\infty(\mathbf{k})}$ is negligible against the paper's
own $\approx 128$-bit post-Grover target (Section~\ref{sec:intro}) only if
$H_\infty(\mathbf{k}) \gtrsim 128$, i.e.\ $K \geq \lceil 128/19.16\rceil = 7$.
$K=1$--$6$ ($H_\infty \approx 19$--$115$ bits) are \emph{not} admissible
under this target: for $K=1$ in particular, $q_F\approx 2^{20}$ offline
HMAC evaluations make the guessing term $\approx 1$, exactly the
brute-force attack Remark~\ref{rem:cvf8} describes. The paper's default
$K=10$ ($H_\infty\approx196$ bits) exceeds this minimum with margin.
Implementations and future revisions of this specification MUST enforce
$K\geq7$ (with the current $|\mathcal{P}|$) and SHOULD retain $K=10$ or
larger; $K<7$ should be treated as a distinct, lower security level, not
as "IND-CPA secure" in the sense of Theorem~\ref{thm:ind-cpa}.
\end{remark}

\begin{remark}[HMAC key-length handling]
\label{rem:hmac-keylen}
$\mathsf{kb}$ has length $5K$ bytes. HMAC-SHA256 hashes any key longer than
its 64-byte block size down to 32 bytes before use, and zero-pads any
shorter key up to 64 bytes; neither transformation is addressed by the
derivation in Section~\ref{sec:construction}. For $K\leq12$ ($5K\leq60$
bytes), $\mathsf{kb}$ is used directly, zero-padded to 64 bytes, and
Theorem~\ref{thm:ind-cpa} applies to $\mathsf{kb}$ as defined. For $K\geq13$
($5K\geq65$ bytes, outside the paper's default), the effective HMAC key
becomes $\mathrm{SHA256}(\mathsf{kb})$: since $\mathsf{kb}\mapsto
\mathrm{SHA256}(\mathsf{kb})$ is collision-resistant and its domain here
($2^{H_\infty(\mathbf{k})}$ possible $\mathsf{kb}$ values) is far smaller
than $2^{256}$, this map is injective with overwhelming probability and
does not reduce $H_\infty(\mathbf{k})$ under the PRF-under-$D$ assumption
of Remark~\ref{rem:prf-d}; Theorem~\ref{thm:ind-cpa} continues to apply
with $\mathsf{kb}$ read as $\mathrm{SHA256}(\mathsf{key\_bytes}(\mathbf{k}))$
in that regime. Either way, the theorem's guessing term is unaffected,
since an adversary can still evaluate the same key-length-handling
transformation offline for each guess at unit cost.
\end{remark}

\begin{remark}[Recommendation: deriving a uniform subkey removes the non-standard assumption]
\label{rem:cvf8-kdf}
Remark~\ref{rem:prf-d} left Theorem~\ref{thm:ind-cpa} resting on
$\mathrm{Adv}^{\mathrm{PRF}}_{\mathrm{HMAC\text{-}SHA256}}(\mathcal{B}_1)$
being negligible \emph{for the specific structured distribution $D$}, not
for a uniform key --- a strictly weaker and non-standard assumption
relative to the literature's usual uniform-key HMAC-SHA256 PRF conjecture.
This residual assumption can be removed entirely, restoring the standard
uniform-key PRF game, by deriving a uniformly-distributed HMAC subkey from
$\mathsf{kb}$ via a KDF before use, e.g.\
$\mathsf{sk} = \mathrm{HKDF\text{-}SHA256}(\mathsf{salt}, \mathsf{kb}, \texttt{"napqes-hmac-key"})$,
and keying every $\mathrm{Derive}_d$ call with $\mathsf{sk}$ in place of
$\mathsf{kb}$. Modelling the KDF as a random oracle (standard for HKDF
extract-then-expand), $\mathsf{sk}$ is statistically close to uniform over
$\{0,1\}^{256}$ regardless of $D$'s structure, so
$\mathrm{Adv}^{\mathrm{PRF}}_{\mathrm{HMAC\text{-}SHA256}}(\mathcal{B}_1)$
in a restated Theorem~\ref{thm:ind-cpa} would then legitimately mean the
standard, uniform-key advantage, at the cost of one extra HKDF evaluation
per key and a wire-format change (this KDF step is not present in the
current reference implementations \texttt{napqes.py}/\texttt{rust/src/lib.rs}/\texttt{C/napqes.c}).
We record this as the preferred long-term fix but have \emph{not} shipped
it as part of the CVF8 response, since it changes the wire format and key
schedule and is out of scope for a proof-only correction; until it lands,
Theorem~\ref{thm:ind-cpa} as restated above (with the explicit
$q_F\cdot2^{-H_\infty(\mathbf{k})}$ term, Remark~\ref{rem:prf-d}'s
PRF-under-$D$ caveat, and Remark~\ref{rem:min-K}'s $K\geq7$ floor) is the
accurate statement of what is actually proved.
\end{remark}

\begin{remark}[Erratum: HKDF-over-$\mathsf{kb}$ does not, by itself, close CVF13]
\label{rem:cvf8-kdf-erratum}
The paragraph above additionally claimed that deriving
$\mathsf{sk}=\mathrm{HKDF}(\mathsf{kb})$ resolves the reduction-simulation
gap of CVF13 (Remark~\ref{rem:cvf13}) ``once $\mathsf{sk}$ is independent
of $\mathbf{k}$.'' This is imprecise and, on reflection, incorrect as a
fix for CVF13 specifically: $\mathsf{sk}=\mathrm{HKDF}(\mathsf{kb})$ is a
\emph{deterministic function} of $\mathbf{k}$ (via
$\mathsf{kb}=\mathsf{key\_bytes}(\mathbf{k})$), so it is not
statistically or computationally independent of $\mathbf{k}$ in the sense
the simulation argument needs --- a reduction that already knows
$\mathbf{k}$ (required to run the arithmetic layer) can compute
$\mathsf{kb}$ and thus $\mathsf{sk}$ itself, so forwarding
domain-derivation calls to an external oracle keyed by
$\mathsf{sk}=\mathrm{HKDF}(\mathsf{kb})$ would be vacuous: the reduction
never needs the oracle in the first place. The leftover-hash-style
argument sketched above only shows $\mathsf{sk}$'s \emph{marginal}
distribution is close to uniform, which is the right notion for CVF8
(key-entropy) but the wrong notion for CVF13 (simulatability), which
requires $\mathsf{sk}$ to be sampled \emph{independently of} $\mathbf{k}$,
not merely to look uniform when $\mathbf{k}$ is also known. Closing CVF13
therefore requires $\mathsf{KeyGen}$ to sample $\mathsf{sk}$ via a fresh,
independent CSPRNG draw --- never as any deterministic function of
$\mathbf{k}$ or $\mathsf{kb}$, HKDF included. Section~\ref{sec:v8} below
specifies and Section~\ref{sec:implementation} implements exactly this:
$\mathsf{KeyGen}_{v8}$ samples $\mathbf{k}$ and $\mathsf{sk}$ independently,
closing both CVF8 (Section~\ref{sec:v8}, uniform $\mathsf{sk}$) and CVF13
(Section~\ref{sec:v8}, genuinely simulatable reductions) simultaneously,
for callers who adopt the v8 schedule.
\end{remark}

\subsection{INT-CTXT: Ciphertext Integrity}

\begin{remark}[Audit finding CVF10: the bound omitted the ideal-world tag-guessing term and was false for $q=0$]
\label{rem:cvf10}
The original statement of Theorem~\ref{thm:int-ctxt} bounded
$\mathrm{Adv}^{\mathrm{INT\text{-}CTXT}}_{\mathrm{NAPQES}}(\mathcal{A})$ by
$\mathrm{Adv}^{\mathrm{PRF}}_{\mathrm{HMAC\text{-}SHA256}}(\mathcal{B}_2) + q^2/2^{128}$
alone. This is unsound: the reduction $\mathcal{B}_2$'s advantage is, by
definition, $\Pr[\text{forge}\mid\text{real}] - \Pr[\text{forge}\mid\text{random}]$,
so a correct bound on $\Pr[\text{forge}\mid\text{real}]$ must carry the
ideal-world term $\Pr[\text{forge}\mid\text{random}]$ forward, not discard
it --- yet even against a \emph{truly random} function, an adversary that
submits a candidate tag $T^*$ for a fresh (never-queried) input succeeds
with probability exactly $2^{-256}$ per attempt, purely by guessing. The
original statement's omission of this term made the bound literally false
at $q=0$: an adversary making \emph{no} encryption queries and submitting
a single uniformly random 256-bit tag as its forgery still wins the
INT-CTXT game with probability $2^{-256} > 0$, strictly exceeding the
original claimed bound of $0$ at $q=0$. The original proof also never
stated $\mathcal{B}_2$'s distinguishing rule, so no advantage accounting
connected $\mathcal{A}$'s forgery to $\mathcal{B}_2$'s output at all. The
restated theorem below adds the missing $q_v/2^{256}$ term, where $q_v$ is
the number of forgery-submission (verification) attempts $\mathcal{A}$
makes, and the proof states $\mathcal{B}_2$'s decision rule explicitly
(guess ``real'' iff $\mathcal{A}$'s forgery verifies). This omission
propagated into the IND-CCA bound (Theorem~\ref{thm:ind-cca}), which
invokes this theorem with $q_d$ verification queries and must therefore
carry a matching $q_d/2^{256}$ term, added below.
\end{remark}

\begin{remark}[Audit finding CVF11: Case 2's ``nonce collision'' term treated an adversarially-chosen value as a random event]
\label{rem:cvf11}
The restated Theorem~\ref{thm:int-ctxt} above (post-CVF10) still carried a
term $q^2/2^{128}$ into Case 2, justified by the sentence ``the probability
of a nonce collision satisfying $N^*=N_i$ for some $i$ is at most
$q/2^{128}\leq q^2/2^{128}$ by the union bound.'' This is unsound: $N^*$ is
not sampled by any experiment. It is a component of the forgery $c^*$ chosen
by $\mathcal{A}$ itself, adaptively, after observing $N_1,\ldots,N_q$ in the
returned ciphertexts. $\mathcal{A}$ can therefore set $N^*=N_i$ for any $i$
of its choosing with probability $1$; this is not a random event over which a
union bound applies, and no such bound may be used to argue $\Pr[N^*=N_i]$
is small. The error is moreover unnecessary: Case 2's own argument already
disposes of $N^*=N_i$ \emph{deterministically}, via the injectivity of the
domain-\texttt{0x03} tag-input encoding (either $\widetilde{B}^*\neq\widetilde{B}_i$,
reducing to Case 1, or $\widetilde{B}^*=\widetilde{B}_i$, in which case the tag
input is identical to a query-phase input and any $T^*\neq T_i$ is rejected
outright) --- so no probability term is needed for Case 2 at all, and the
$q^2/2^{128}$ term is removed below. The propagation of this term into the
IND-CCA bound (Theorem~\ref{thm:ind-cca}, via ``$q+q_d$ total oracle
queries'') is removed for the same reason; the genuine $q^2/2^{128}$ term
that remains in Theorem~\ref{thm:ind-cca} comes only from the honestly-random
challenge nonce in the IND-CPA hybrid (Lemma~\ref{lem:hiding}'s
$\mathrm{Coll}$ event), not from INT-CTXT.
\end{remark}

\begin{remark}[Audit finding CVF12: stale pre-CVF2 tag-input formula survived in Theorem~\ref{thm:int-ctxt}'s Case~1 and the AAD Binding paragraph]
\label{rem:cvf12}
Theorem~\ref{thm:int-ctxt}'s Case~1 argument and the AAD Binding paragraph
below its proof both wrote the tag input as
$\texttt{0x03}\|\mathrm{be4}(|aad|)\|aad\|N\|\widetilde{B}$ --- the
\emph{pre}-CVF2, AAD-first-then-nonce layout, in which the nonce sits at a
variable byte offset depending on $|aad|$ --- even though
Remark~\ref{rem:domsep} (the CVF2 fix) and the hiding lemma's
``Freshness of the authentication tag'' paragraph (Lemma~\ref{lem:hiding})
already state, correctly, that domain \texttt{0x03} uses the unified
$d\|N\|\mathsf{ctx}$ layout with the nonce at the fixed offset $1..16$,
matching the shipped code (\texttt{napqes.py}'s
\texttt{\_compute\_auth\_tag}, and the equivalent Rust/C routines). This
is a real internal inconsistency: the same construction was described by
two different, mutually incompatible formulas in the same document, and
the stale formula is precisely the one under which the audit's
cross-group collision scenario is possible --- a nonce-first input from
another domain, $N_i\|d\|\mathsf{ctx}$ (21--22 bytes for the
fixed-width-counter domains), can equal an AAD-first tag input
$\texttt{0x03}\|\mathrm{be4}(0)\|N^*$ (21 bytes, the empty-AAD,
$|\widetilde{B}^*|=0$ minimal case) whenever $N_i$ happens to begin with
$\texttt{0x03}\,\texttt{00}\,\texttt{00}\,\texttt{00}\,\texttt{00}$, an
event of probability $\approx 2^{-40}$ per query (or $\approx q/2^{40}$
over $q$ queries by the union bound) --- exactly the term the audit
finding identifies as missing under that (stale) formula. Under the
\emph{actual}, already-unified domain-first layout, this collision cannot
occur at all: Remark~\ref{rem:domsep}'s cross-domain argument (inputs
with different domain bytes differ at byte position~$0$, unconditionally)
already rules out any \texttt{0x03} input colliding with a nonce-first
domain's input, with no probabilistic term of any kind. The fix below
corrects Theorem~\ref{thm:int-ctxt}'s Case~1 and the AAD Binding
paragraph to the formula Remark~\ref{rem:domsep} and the hiding lemma
already use, and grounds Case~1's freshness argument in injectivity (per
Remark~\ref{rem:domsep}) rather than the previously imprecise ``prefixed
by nonces'' phrasing --- matching the audit's own recommended fix for the
domain-first case. We also rechecked Remark~\ref{rem:domsep} for the
audit's separately-flagged claim that a $q/2^{32}$ cross-group term is
``dominated by'' the nonce-collision term $q^2/2^{128}$ (mathematically
false for realistic $q$, as the finding correctly notes); no such claim is
present anywhere in the current Remark~\ref{rem:domsep} --- the CVF2
rewrite already replaced the probabilistic cross-group argument with the
unconditional injectivity argument, so there is no $q/2^{32}$ term left to
(mis)dominate, and no bound in this document ever carried one. No further
fix is required for that sub-claim beyond confirming its absence.
\end{remark}

\begin{theorem}[INT-CTXT]\label{thm:int-ctxt}
  Under the PRF assumption on HMAC-SHA256, NAPQES is INT-CTXT secure.
  Specifically, for any PPT adversary $\mathcal{A}$ making at most $q$
  encryption queries and at most $q_v \geq 1$ forgery-submission
  (verification) attempts, there exists a PRF adversary $\mathcal{B}_2$
  with similar running time such that:
  \[
    \mathrm{Adv}^{\mathrm{INT\text{-}CTXT}}_{\mathrm{NAPQES}}(\mathcal{A})
    \leq
    \mathrm{Adv}^{\mathrm{PRF}}_{\mathrm{HMAC\text{-}SHA256}}(\mathcal{B}_2)
    + \frac{q_v}{2^{256}}.
  \]
  For a single forgery attempt ($q_v = 1$, the minimal INT-CTXT
  experiment), the last term is $2^{-256}$; it is what makes the bound
  correctly nonzero at $q=0$ (Remark~\ref{rem:cvf10}). There is no
  $q$-dependent term: $q$ (the number of encryption queries) affects only
  which nonces $\mathcal{A}$ has seen, and Case 2 below disposes of
  $N^*=N_i$ deterministically rather than probabilistically
  (Remark~\ref{rem:cvf11}). The experiment grants $\mathcal{A}$ a genuine
  encryption oracle and a genuine, \emph{adaptive} decryption/verification
  oracle: $q_v$ counts real-time oracle queries, each answered with the
  true decrypted plaintext or $\bot$ before $\mathcal{A}$ chooses its next
  query, not a batch of final guesses evaluated only once at the end
  (Remark~\ref{rem:cvf21}).
\end{theorem}

\begin{remark}[Audit finding CVF20: Case~2 ignored $aad^*$, and freshness was defined on $c^*$ alone]
\label{rem:cvf20}
An earlier draft of this proof (a) defined freshness as $c^* \notin
\{c_1,\ldots,c_q\}$ --- a condition on the ciphertext alone --- and (b)
split Case~2 into $\widetilde{B}^* \neq \widetilde{B}_i$ versus
$\widetilde{B}^* = \widetilde{B}_i \wedge T^* \neq T_i$, asserting in the
second sub-case that ``the tag input is identical to $x_i$.'' That claim
holds only if additionally $aad^* = aad_i$, which was never assumed: if
$aad^* \neq aad_i$ while $\widetilde{B}^* = \widetilde{B}_i$, then
$x^* \neq x_i$ (Table~\ref{tab:domains}'s length-prefixed
$\mathrm{be4}(|aad|)\|aad$ encoding is injective in $aad$), and this
sub-case must route to the Case~1 fresh-input argument rather than being
dismissed as unable to contribute. The standard INT-CTXT object of
freshness is the pair $(c,aad)$, not $c$ alone: a forgery that replays
$c_i$ verbatim with a new $aad^* \neq aad_i$ is a pair never returned by
the encryption oracle, hence a legitimate forgery attempt, and (a) silently
excluded it rather than proving it infeasible. The fix below (i) restates
freshness over $(c^*,aad^*)$ and (ii) re-splits Case~2 on $x^*=x_i$ versus
$x^*\neq x_i$, which subsumes and correctly routes the AAD-replay case.
\end{remark}

\begin{remark}[Audit finding CVF21: Theorem~\ref{thm:int-ctxt} must be proved in the adaptive, multi-verification-query oracle regime the IND-CCA proof actually invokes it in]
\label{rem:cvf21}
Theorem~\ref{thm:int-ctxt}'s statement described $q_v$ only as a count of
``forgery-submission attempts,'' which reads as a batch of final guesses
evaluated once at the end of the experiment. The IND-CCA proof
(Theorem~\ref{thm:ind-cca}), however, invokes this theorem to bound
$\Pr[D]$ for a forger $\mathcal{A}'$ that holds a genuine, real-time
decryption oracle $\hat{\mathcal{D}}$ and answers each of $\mathcal{A}$'s
$q_d$ decryption queries adaptively, with $\mathcal{A}$ observing each
$\bot$/plaintext answer before choosing its next query --- a strictly
stronger, Bellare--Namprempre-style multi-query oracle regime that the
original statement, read literally, never established. Two further gaps
compounded this, per the finding: (i) an earlier draft's Case~2 carried a
$q^2/2^{128}$ term into the IND-CCA bound via ``$q+q_d$ total oracle
queries'' (Remark~\ref{rem:cvf11}), substituting $q+q_d$ for the
encryption-query parameter $q$ of a term that was, in any case, a
collision bound over \emph{encryption-generated} nonces --- decryption
queries generate no nonces, so that substitution had no basis in the
theorem being invoked, independently of the term's own soundness; and
(ii) the forger $\mathcal{A}'$ constructed in Theorem~\ref{thm:ind-cca}'s
proof must itself produce the IND-CCA challenge ciphertext $c^*$ (by
sampling $b$ and querying its own encryption oracle on $m_b$) to run the
post-challenge phase of $\mathcal{A}$ at all --- $q+1$ encryption queries,
not $q$ --- yet an earlier draft's simulation description omitted this
challenge-generation step entirely while asserting, rather than showing,
that ``$\mathcal{A}'$'s simulation is perfect.''

\textbf{Resolution.} We resolve this by proving the bound holds in
exactly the adaptive oracle regime required, rather than restating it as
a weaker, single-shot game:
\begin{itemize}
  \item \emph{Adaptivity does not change the Case~1/Case~2 argument.}
        Fix any sequence of $q_v$ adaptive decryption/verification queries,
        each answered in real time before the next is chosen. Case~2
        (above) is resolved by injectivity alone --- for a query with
        $N=N_i$ for some encryption-query $i$, whether $x=x_i$ or
        $x\neq x_i$ is fully determined by the query's own bytes,
        independent of any other query's answer or ordering, so
        adaptivity contributes nothing to Case~2 regardless of how many
        queries have already been answered. Case~1 (a query with a fresh
        $N$) reduces the query's success probability, in the ideal
        world, to $f(x)$ for an $x$ that --- by the same injectivity
        argument --- has never been queried to $f$ before, \emph{even
        counting every previous decryption/verification query's own $x$
        as part of $f$'s query history}: lazy sampling of the random
        function $f$ gives every not-yet-queried input an independent
        uniform $256$-bit value regardless of what other, distinct inputs
        have already been resolved, so knowledge of every previous
        query's outcome (accept or $\bot$) reveals nothing about $f$ at a
        new, still-fresh point. The per-query success probability is
        therefore exactly $2^{-256}$ regardless of adaptivity, position in
        the query sequence, or $\mathcal{A}$'s strategy for choosing later
        queries from earlier answers, and the union bound
        $\Pr[\text{forge in Case 1 across all } q_v \text{ queries}] \leq
        q_v/2^{256}$ therefore holds for the fully adaptive, real-time
        oracle exactly as it did for a final batch of guesses --- the
        \emph{bound} of Theorem~\ref{thm:int-ctxt} is unchanged, but it is
        now established for the stronger regime the IND-CCA proof
        requires, closing this half of the finding.
  \item \emph{No substitution of $q+q_d$ into a nonce-collision parameter
        is needed}, because none exists to substitute into:
        Remark~\ref{rem:cvf11} already removed Case~2's $q^2/2^{128}$ term
        from Theorem~\ref{thm:int-ctxt} entirely (it traced to a bogus
        argument over the adversarially-chosen $N^*$, not to a genuine
        collision over encryption-generated nonces), so
        Theorem~\ref{thm:int-ctxt}'s bound has no $q$-dependent term at
        all for the IND-CCA proof to (mis)parametrise by $q+q_d$. This was
        gap~(i) above; it is moot under the current bound rather than
        requiring a corrected substitution.
  \item \emph{The forger's own challenge query, and the perfect-simulation
        claim, are made explicit} in the revised construction of
        $\mathcal{A}'$ below (Theorem~\ref{thm:ind-cca}'s proof):
        $\mathcal{A}'$ is shown, step by step, to sample $b$ itself and
        query its own encryption oracle on $m_b$ as its $(q\!+\!1)$-th
        encryption query, to record the resulting $c^*$ for later
        freshness bookkeeping (Remark~\ref{rem:cvf23}), and to forward
        every one of $\mathcal{A}$'s decryption queries --- pre- \emph{and}
        post-challenge (Remark~\ref{rem:cvf22}) --- to its own real
        decryption oracle; perfection of the simulation is then derived
        from the fact that every oracle $\mathcal{A}$ interacts with is a
        real oracle keyed by the same hidden key throughout, rather than
        merely asserted.
\end{itemize}
\end{remark}

\begin{proof}
The INT-CTXT experiment requires $\mathcal{A}$ to produce a \emph{fresh}
valid ciphertext: a pair $(c^*, aad^*)$ with
$(c^*, aad^*) \notin \{(c_1,aad_1), \ldots, (c_q,aad_q)\}$ (the outputs of
the $q$ encryption oracle calls, paired with their queried AAD;
Remark~\ref{rem:cvf20}) such that
$\mathrm{NAPQES}.\mathsf{Dec}(\mathbf{k}, c^*, aad^*) \neq \bot$. In
particular $\mathcal{A}$ may submit $c^* = c_i$ for some $i$ provided
$aad^* \neq aad_i$ --- a replay-with-new-AAD attempt, addressed in Case~2
below.

Write $c^* = (N^*, \widetilde{B}^*, T^*)$.  For $\mathsf{Dec}$ to accept,
$T^*$ must satisfy (domain-first layout, Remark~\ref{rem:domsep};
corrected here per audit finding CVF12, Remark~\ref{rem:cvf12} --- an
earlier draft of this proof used the pre-CVF2, AAD-first-then-nonce
formula):
\[
  T^* = \mathrm{HMAC}\!\left(\mathsf{kb},\
        \texttt{0x03} \| N^* \| \mathrm{be4}(|aad^*|) \| aad^* \| \widetilde{B}^*\right).
\]
Let $x^* = \texttt{0x03} \| N^* \| \mathrm{be4}(|aad^*|) \| aad^* \| \widetilde{B}^*$.

\paragraph{Case 1: $N^*$ was not used in any encryption query.}
Every domain-\texttt{0x03} tag input from the query phase has the form
$x_i = \texttt{0x03} \| N_i \| \mathrm{be4}(|aad_i|) \| aad_i \| \widetilde{B}_i$,
with $N_i$ occupying the same fixed byte offset $1..16$ that $N^*$
occupies in $x^*$ (Remark~\ref{rem:domsep}: every domain-\texttt{0x03}
input shares the uniform $d\|N\|\mathsf{ctx}$ shape). Since $N_i \neq N^*$
for every query $i$ (Case~1 hypothesis), $x_i \neq x^*$ unconditionally at
bytes $1..16$, regardless of $|aad_i|$, $|aad^*|$, or $\widetilde{B}_i$ ---
freshness follows from injectivity of the encoding, not from an assumed
nonce prefix, and no probabilistic cross-group collision term (of the
kind audit finding CVF12 checked for; see Remark~\ref{rem:cvf12}) is
needed or applies: Remark~\ref{rem:domsep}'s cross-domain distinctness
additionally rules out $x^*$ colliding with any other domain's input. So
$x^*$ was never submitted to $\mathrm{HMAC}(\mathsf{kb}, \cdot)$ during
the query phase.  We construct
$\mathcal{B}_2$, given oracle access to a function $f$ that is either
$\mathrm{HMAC}(\mathsf{kb}, \cdot)$ (the \emph{real} world) or a uniformly
random function $\{0,1\}^* \to \{0,1\}^{256}$ (the \emph{ideal} world):
$\mathcal{B}_2$ simulates the INT-CTXT experiment for $\mathcal{A}$,
forwarding every HMAC computation to $f$, and evaluates each of
$\mathcal{A}$'s (at most $q_v$) forgery-submission attempts by computing
the expected tag via $f$. This description shares, for the same reason,
the simulation gap identified in Remark~\ref{rem:cvf13}: answering
$\mathcal{A}$'s $q$ encryption queries requires running
$\mathrm{NAPQES}.\mathsf{Enc}(\mathbf{k},\cdot)$, which needs $\mathbf{k}$
itself for its arithmetic layer, not merely oracle access to $f$; see that
remark for the resolution. \emph{Decision rule:} $\mathcal{B}_2$ outputs
``real'' iff at least one of $\mathcal{A}$'s forgeries verifies (its
candidate tag matches $f(x^*)$ for the corresponding tag input $x^*$);
otherwise it outputs ``random''. By this construction,
\[
  \mathrm{Adv}^{\mathrm{PRF}}_{\mathrm{HMAC\text{-}SHA256}}(\mathcal{B}_2)
  =
  \Pr[\mathcal{B}_2 \to \text{``real''} \mid f = \mathrm{HMAC}(\mathsf{kb},\cdot)]
  -
  \Pr[\mathcal{B}_2 \to \text{``real''} \mid f \text{ random}]
  =
  \Pr[\text{forge}\mid\text{real}] - \Pr[\text{forge}\mid\text{random}].
\]
In the real world this probability is exactly $\mathcal{A}$'s Case-1
forgery probability. In the ideal world, $x^*$ was never queried, so
$f(x^*)$ is a uniform 256-bit string independent of $\mathcal{A}$'s view;
each of $\mathcal{A}$'s $q_v$ forgery attempts therefore succeeds only by
blind guessing, with probability exactly $2^{-256}$ per attempt, so by the
union bound $\Pr[\text{forge}\mid\text{random}] \leq q_v/2^{256}$.
Rearranging the advantage equation,
\[
  \Pr[\text{forge in Case 1}]
  \;\leq\;
  \mathrm{Adv}^{\mathrm{PRF}}_{\mathrm{HMAC\text{-}SHA256}}(\mathcal{B}_2)
  + \frac{q_v}{2^{256}}.
\]
This ideal-world term is exactly what the original proof omitted
(Remark~\ref{rem:cvf10}): it is what makes the bound correct at $q=0$,
since an adversary submitting one uninformed random tag guess still
forges with probability $2^{-256}$ even against a perfect random
function.

\paragraph{Case 2: $N^* = N_i$ for some query $i$.}
Fix any query $i$ with $N_i = N^*$, and recall
$x_i = \texttt{0x03} \| N_i \| \mathrm{be4}(|aad_i|) \| aad_i \| \widetilde{B}_i$.
The correct split (Remark~\ref{rem:cvf20}) is on $x^* = x_i$ versus
$x^* \neq x_i$, not on $\widetilde{B}^*$ alone:
\begin{itemize}
  \item \emph{$x^* \neq x_i$ for every query $i$ with $N_i = N^*$.} This
        covers $\widetilde{B}^* \neq \widetilde{B}_i$ \emph{and} the
        case $\widetilde{B}^* = \widetilde{B}_i$ with $aad^* \neq aad_i$
        (by injectivity of $\mathrm{be4}(|aad|)\|aad$ in $aad$, holding
        $N_i=N^*$ and $\widetilde{B}^*=\widetilde{B}_i$ fixed) --- in
        particular the replay-with-new-AAD attempt
        $(c^*,aad^*)=(c_i,aad^*)$, $aad^*\neq aad_i$, falls here. In every
        such case $x^*$ was never submitted to
        $\mathrm{HMAC}(\mathsf{kb},\cdot)$ during the query phase (it
        differs from every $x_j$: from $x_i$ by the argument just given,
        and from every $x_j$ with $N_j \neq N^*$ by the Case~1 argument
        applied to that $j$), so the Case~1 reduction and bound apply
        verbatim to this forgery attempt.
  \item \emph{$x^* = x_i$ for some query $i$ with $N_i = N^*$.} By
        injectivity of the tag-input encoding (Remark~\ref{rem:domsep}),
        $x^*=x_i$ forces $aad^*=aad_i$ and $\widetilde{B}^*=\widetilde{B}_i$
        jointly. Two sub-cases on the tag itself: if $T^* = T_i$, then
        $c^*=(N^*,\widetilde{B}^*,T^*)=c_i$ and $aad^*=aad_i$, so
        $(c^*,aad^*)=(c_i,aad_i)$ --- excluded by the freshness
        requirement (Remark~\ref{rem:cvf20}), hence not a valid forgery
        attempt. If instead $T^* \neq T_i$, decryption recomputes the
        expected tag as $\mathrm{HMAC}(\mathsf{kb},x^*) =
        \mathrm{HMAC}(\mathsf{kb},x_i) = T_i \neq T^*$ and rejects
        outright, regardless of freshness. Either way this branch
        contributes no successful, freshness-respecting forgery.
\end{itemize}
Both branches are exhaustive and are resolved \emph{deterministically} by
injectivity, not by any probabilistic argument over $N^*$: $N^*$ is a
value $\mathcal{A}$ chooses as part of its forgery $c^*$, adaptively and
after observing $N_1,\ldots,N_q$, not a value sampled by the experiment,
so no probability term (in particular, no $q^2/2^{128}$ nonce-collision
term) is associated with Case~2 (Remark~\ref{rem:cvf11}). The re-split
above additionally shows the replay-with-new-AAD attempt is not silently
excluded but is proved infeasible beyond the Case~1 bound: it always
falls in the first branch, which reduces to Case~1. Case~2 therefore
contributes no additional term beyond Case~1's, and combining both cases:
\[
  \mathrm{Adv}^{\mathrm{INT\text{-}CTXT}}_{\mathrm{NAPQES}}(\mathcal{A})
  \;\leq\;
  \mathrm{Adv}^{\mathrm{PRF}}_{\mathrm{HMAC\text{-}SHA256}}(\mathcal{B}_2)
  + \frac{q_v}{2^{256}}. \qedhere
\]
\end{proof}

\subsection{AAD Binding}

The AAD is included in the tag computation as (domain-first layout,
Remark~\ref{rem:domsep}; corrected here per audit finding CVF12,
Remark~\ref{rem:cvf12}):
$T = \mathrm{HMAC}(\mathsf{kb},\ \texttt{0x03} \| N \| \mathrm{be4}(|A|) \| A \| B)$.
The 4-byte length prefix prevents length-extension mixing between different
AAD lengths.  Any single-bit change to the AAD produces an unpredictable
change to the expected tag, with advantage bounded by the HMAC security.

\subsection{IND-CCA Security}
\label{sec:ind-cca}

\begin{remark}[Audit finding CVF14: the $H_0\to H_1$ ``identical-until-bad'' step is false, since $H_1$ also diverges on non-fresh replays]
\label{rem:cvf14}
An earlier draft of the proof below defined $H_1$ by stubbing
$\mathcal{D}$ to \emph{always} return $\bot$, while defining the bad event
$D$ as a \emph{fresh} ($c \notin \{c_1,\ldots,c_q\}$) valid submission, and
then asserted ``$H_0$ and $H_1$ are identical in every execution that does
not trigger $D$.'' That assertion is false as those two definitions stand:
in $H_0$, $\mathcal{A}$ may legally query $\mathcal{D}$ on a
\emph{replayed} encryption-oracle output $(c_i, aad_i)$ and receive the
correct plaintext $m_i$ (by correctness, Definition~\ref{def:correctness}),
whereas the stubbed $H_1$ returns $\bot$ on that same query --- yet this
query is not fresh, so $D$ does not occur. An adversary that simply
queries $\mathcal{D}(c_1, aad_1)$ after one encryption-oracle query
distinguishes $H_0$ from $H_1$ with probability $1$ without ever
triggering $D$, so the fundamental-lemma step
$|\Pr[H_0]-\Pr[H_1]|\leq\Pr[D]$ was not established as written. This is a
distinct defect from CVF10--CVF13, which concern gaps or omitted terms
inside the INT-CTXT and IND-CCA advantage bounds themselves; CVF14
concerns whether the hybrid argument used to \emph{reduce} IND-CCA to
INT-CTXT plus IND-CPA is even correctly stated, independent of those
bounds' content. The fix, applied below, is the standard corrected hybrid:
redefine $H_1$ so that $\mathcal{D}$ answers a replayed $(c_i, aad_i)$
pair with its recorded plaintext $m_i$ from a table populated by the
encryption oracle, and returns $\bot$ only for ciphertexts not in that
table (i.e., only for fresh queries). Under that definition $H_0$ and
$H_1$ agree exactly on replayed queries (both return $m_i$, by
correctness) and can only disagree on a fresh query, which is exactly
what $D$ captures, so the identical-until-bad step now goes through.
\end{remark}

\begin{remark}[Audit finding CVF22: the bad event $D$ was restricted to the query phase, leaving post-challenge forgeries uncovered]
\label{rem:cvf22}
Game $H_0$ below explicitly permits $\mathcal{A}$ to call the decryption
oracle after receiving the challenge ciphertext, yet an earlier draft
defined the bad event $D$ as a fresh valid submission occurring only
``during the query phase'' (i.e., before the challenge). A fresh valid
ciphertext submitted \emph{after} the challenge makes $H_0$ and $H_1$
diverge (by exactly the same case analysis as a pre-challenge query) but
was not counted by that restricted $D$, so the identical-until-bad bound
$|\Pr[H_0]-\Pr[H_1]|\leq\Pr[D]$ again failed to cover every
distinguishing execution --- the same defect family as CVF14, applied to
the phase of the query rather than to replays. \textbf{Fix:} $D$ below is
defined over decryption queries in \emph{both} phases (any fresh valid
$(c,aad)$ submitted to $\mathcal{D}$ at any point in the experiment), and
the INT-CTXT forger $\mathcal{A}'$ constructed in the proof forwards
\emph{every} one of $\mathcal{A}$'s decryption queries --- pre- and
post-challenge alike --- to its own real decryption oracle, so the
reduction bounding $\Pr[D]$ runs across both phases exactly as $D$ is now
defined.
\end{remark}

\begin{remark}[Audit finding CVF23: freshness-set mismatch, and $c$-only freshness silently weakened the IND-CCA notion]
\label{rem:cvf23}
Two related defects. First, an earlier draft's bad event $D$ used the
freshness set $\{c_1,\ldots,c_q\}$ (the adversary's own $q$ encryption
queries), but from the INT-CTXT challenger's point of view the challenge
ciphertext $c^*$ is itself an $\hat{\mathcal{E}}$ output once the forger
$\mathcal{A}'$ generates it (Remark~\ref{rem:cvf21}), so the correct
freshness set for $\mathcal{A}'$'s win condition is
$\{c_1,\ldots,c_q\}\cup\{c^*\}$; the two sets were never reconciled.
Second, and independently, an earlier draft's Game~$H_0$ forbade
$\mathcal{A}$ from querying $\mathcal{D}$ on \emph{any} $c=c^*$ (for every
$aad$), whereas the standard AEAD IND-CCA game forbids only the specific
pair $(c^*,aad^*)$ and \emph{permits} $\mathcal{D}(c^*,aad\neq aad^*)$ ---
a legitimate, potentially-winning forgery attempt in the standard game
that the $c$-only restriction silently excluded, consistent with the
$c$-only (rather than $(c,aad)$-pair) freshness already corrected for
Theorem~\ref{thm:int-ctxt} by audit finding CVF20
(Remark~\ref{rem:cvf20}). Left uncorrected, the composed result would
establish a strictly weaker-than-standard IND-CCA notion --- one in which
$\mathcal{A}$ is never even permitted to \emph{attempt} the
AAD-substitution forgery, rather than one in which that forgery is proved
to fail. \textbf{Fix:} the table $T$ used by $D$/$D'$ below is populated
to include the challenge tuple $(c^*,aad^*,m_b)$ once the challenge is
issued, so $\mathcal{A}'$'s freshness set is exactly
$\{(c_1,aad_1),\ldots,(c_q,aad_q)\}\cup\{(c^*,aad^*)\}$ --- the INT-CTXT
challenger's own complete output set, reconciling the mismatch --- and
Game~$H_0$ now forbids only the pair $(c^*,aad^*)$, explicitly permitting
$\mathcal{D}(c^*,aad\neq aad^*)$, so the proven notion is the standard
$(c,aad)$-pair-freshness AEAD IND-CCA game rather than the $c$-only
weakening.
\end{remark}

\begin{theorem}[IND-CCA]\label{thm:ind-cca}
  Under the PRF assumption on HMAC-SHA256, NAPQES is IND-CCA secure.
  Specifically, for any PPT adversary $\mathcal{A}$ making at most $q$
  encryption queries and $q_d$ decryption queries there exist PRF adversaries
  $\mathcal{B}_1$, $\mathcal{B}_2$ with similar running time such that:
  \[
    \mathrm{Adv}^{\mathrm{IND\text{-}CCA}}_{\mathrm{NAPQES}}(\mathcal{A})
    \leq
    \mathrm{Adv}^{\mathrm{PRF}}_{\mathrm{HMAC\text{-}SHA256}}(\mathcal{B}_1)
    +
    \mathrm{Adv}^{\mathrm{PRF}}_{\mathrm{HMAC\text{-}SHA256}}(\mathcal{B}_2)
    + \frac{q^2}{2^{128}}
    + \frac{q_d}{2^{256}}.
  \]
  The $q_d/2^{256}$ term is the ideal-world forgery term of
  Theorem~\ref{thm:int-ctxt} (Remark~\ref{rem:cvf10}), carried through
  with $q_v = q_d$ since the INT-CTXT forger $\mathcal{A}'$ constructed
  below makes exactly $q_d$ verification (decryption-oracle) queries,
  counted across \emph{both} the pre- and post-challenge phases
  (Remark~\ref{rem:cvf22}). Freshness for $D$ and for the INT-CTXT
  reduction is measured over the pair $(c,aad)$ against the set
  $\{(c_1,aad_1),\ldots,(c_q,aad_q)\}\cup\{(c^*,aad^*)\}$
  (Remark~\ref{rem:cvf23}). The forger $\mathcal{A}'$ constructed below
  additionally uses $q+1$ total encryption-oracle queries (its $q$
  forwarded queries plus one to produce $c^*$ itself) when invoking
  Theorem~\ref{thm:int-ctxt} --- which carries no
  encryption-query-count-dependent term (Remark~\ref{rem:cvf11}) and is
  proved in the adaptive dual-oracle regime this invocation requires
  (Remark~\ref{rem:cvf21}) --- so this does not change the stated bound.
  The $q^2/2^{128}$ term is the genuine nonce-collision bound from the
  IND-CPA hybrid (Lemma~\ref{lem:hiding}'s $\mathrm{Coll}$ event over the
  $q$ honestly-random encryption-oracle nonces plus the challenge nonce);
  it is \emph{not} inherited from Theorem~\ref{thm:int-ctxt}, which has no
  $q$-dependent term (Remark~\ref{rem:cvf11}).
  (As noted in Remark~\ref{rem:cvf14}, the hybrid game $H_1$ used in the
  proof below is defined with a table-backed decryption oracle rather
  than one stubbed unconditionally to $\bot$, so that the
  identical-until-bad step is sound.)
\end{theorem}

\begin{proof}
We apply the composition theorem of Bellare and
Namprempre~\cite{bellare2000ae} (B\&N~2000, Theorem~3): a scheme that
is simultaneously IND-CPA and INT-CTXT is IND-CCA.  We make the argument
explicit via two hybrid games. (As noted in Remark~\ref{rem:cvf13}, the
$\mathrm{Adv}^{\mathrm{PRF}}(\mathcal{B}_1)$ and
$\mathrm{Adv}^{\mathrm{PRF}}(\mathcal{B}_2)$ terms this composition
invokes below inherit that remark's simulation-gap caveat.)

\paragraph{Game $H_0$ (Real IND-CCA).}
The standard IND-CCA experiment: $\mathcal{A}$ has access to an encryption
oracle $\mathcal{E}$ and a decryption oracle $\mathcal{D}$, and may query
$\mathcal{D}$ both before and after receiving the challenge ciphertext
$c^*$ (produced by one call to $\mathcal{E}$ on adversarially-chosen
$(A^*,m_b)$). The \emph{only} restriction is on the exact pair:
$\mathcal{A}$ may not query $\mathcal{D}(c^*, A^*)$, but \emph{may} query
$\mathcal{D}(c^*, aad)$ for any $aad \neq A^*$ (Remark~\ref{rem:cvf23}) ---
the standard AEAD $(c,aad)$-pair freshness notion, not the strictly
weaker $c$-only restriction (``$\mathcal{D}$ on any $c \neq c^*$'') an
earlier draft imposed.

\paragraph{Game $H_1$ (Table-backed decryption oracle).}
Identical to $H_0$, except $\mathcal{D}$ is replaced by $\mathcal{D}'$,
which consults a table $T$ populated by every $(c, aad, m)$ triple
$\mathcal{E}$ has produced so far --- including, once issued, the
challenge triple $(c^*, A^*, m_b)$ itself (Remark~\ref{rem:cvf23}: from
the standpoint of freshness, the challenge is an $\mathcal{E}$ output like
any other): on query $(c, aad)$, $\mathcal{D}'$ returns the recorded $m$
if $(c, aad)$ matches some entry in $T$, and returns $\bot$ otherwise ---
i.e., $\bot$ only on pairs not yet in $T$ (per Remark~\ref{rem:cvf14}; an
earlier draft instead stubbed $\mathcal{D}$ to return $\bot$
unconditionally, which is corrected here). The one query the game
forbids outright, $\mathcal{D}(c^*,A^*)$, would --- were it permitted ---
trivially hit the table entry for $(c^*,A^*,m_b)$; since $\mathcal{A}$ is
never allowed to make it, $\mathcal{D}'$ never needs to disclose $m_b$ to
answer any query it actually receives.

\paragraph{Transition $H_0 \to H_1$.}
Let $D$ be the event that $\mathcal{A}$ submits, \emph{at any point in the
experiment --- before or after the challenge alike}
(Remark~\ref{rem:cvf22}), a \emph{fresh valid} ciphertext to
$\mathcal{D}$—i.e., a pair $(c, aad)$ with $(c, aad) \notin T$ (not equal
to any $(c_i, aad_i)$ recorded from a pre- or post-challenge
encryption-oracle output, nor to $(c^*,A^*)$ itself,
Remark~\ref{rem:cvf23}) such that
$\mathrm{NAPQES}.\mathsf{Dec}(\mathbf{k}, c, aad) \neq \bot$.
We claim $H_0$ and $H_1$ answer every decryption query identically unless
$D$ occurs, in \emph{either} phase. Consider any query $(c, aad)$
$\mathcal{A}$ makes to the decryption oracle, pre- or post-challenge: (a)
if $(c, aad) \in T$, i.e. it replays some earlier $\mathcal{E}$ output
(a replay of $(c^*,A^*)$ itself cannot occur, since $H_0$ forbids that
exact query), then by correctness (Definition~\ref{def:correctness})
$H_0$'s real decryption returns exactly the recorded plaintext, which is
precisely what $H_1$'s table lookup returns --- the two games agree on
every replayed query, closing exactly the gap identified in
Remark~\ref{rem:cvf14}; (b) if $(c, aad) \notin T$ and the ciphertext is
invalid, both games return $\bot$; (c) if $(c, aad) \notin T$ and the
ciphertext is valid, $H_0$ returns the (non-$\bot$) decrypted plaintext
while $H_1$ returns $\bot$ --- this is exactly event $D$. This case
analysis is identical in the pre-challenge and post-challenge phases ---
nothing about it depends on whether the challenge has been issued yet,
only on whether $(c,aad)\in T$ --- so restricting $D$ to a single phase,
as an earlier draft did, was an unjustified narrowing
(Remark~\ref{rem:cvf22}) rather than a consequence of the case analysis
itself. Hence $H_0$ and $H_1$ are identical in every execution that does
not trigger $D$ in \emph{either} phase, and the fundamental lemma of
game-hopping gives:
\[
  \bigl|\Pr[\mathcal{A}\text{ wins }H_0] - \Pr[\mathcal{A}\text{ wins }H_1]\bigr|
  \;\leq\; \Pr[D].
\]

We bound $\Pr[D]$ by constructing an INT-CTXT forger $\mathcal{A}'$ (using
the B\&N formulation that grants the INT-CTXT adversary both an encryption
oracle \emph{and} an adaptive decryption/verification oracle, in the
regime proved sound by Remark~\ref{rem:cvf21}~\cite{bellare2000ae},
Definition~2):
\begin{enumerate}
  \item $\mathcal{A}'$ receives an encryption oracle $\hat{\mathcal{E}}$ and a
        decryption oracle $\hat{\mathcal{D}}$ from the INT-CTXT challenger,
        and initialises an empty table $T$.
  \item $\mathcal{A}'$ runs $\mathcal{A}_1$, forwarding every encryption
        query to $\hat{\mathcal{E}}$ (recording each resulting
        $(c_i, aad_i, m_i)$ triple into $T$) and every (pre-challenge)
        decryption query to $\hat{\mathcal{D}}$: whenever $\hat{\mathcal{D}}$
        returns a non-$\bot$ result on a pair $(c,aad)\notin T$, event $D$
        has occurred and $\mathcal{A}'$ immediately wins the INT-CTXT
        experiment.
  \item When $\mathcal{A}_1$ outputs $(A^*, m_0, m_1, \mathsf{st})$,
        $\mathcal{A}'$ samples $b \xleftarrow{\$} \{0,1\}$ itself and
        queries $\hat{\mathcal{E}}$ on $(A^*, m_b)$ --- its
        $(q\!+\!1)$-th encryption query --- to obtain the challenge
        ciphertext $c^*$, records $(c^*, A^*, m_b)$ into $T$
        (Remark~\ref{rem:cvf23}), and invokes $\mathcal{A}_2(\mathsf{st},
        c^*)$: the challenge phase is thus an explicit, counted step of
        the construction, not an unstated assumption
        (Remark~\ref{rem:cvf21}, gap~(ii)).
  \item $\mathcal{A}'$ continues running $\mathcal{A}_2$, forwarding
        every \emph{post-challenge} decryption query to $\hat{\mathcal{D}}$
        exactly as in step~2 --- $\mathcal{A}_2$ may query any
        $(c,aad)\neq(c^*,A^*)$, per Game~$H_0$'s restriction
        (Remark~\ref{rem:cvf23}) --- and again declares event $D$ (and
        wins) on the first non-$\bot$ result for a pair not in $T$; this
        step is what makes the reduction cover the post-challenge phase,
        which an earlier draft's construction omitted
        (Remark~\ref{rem:cvf22}).
  \item If $\mathcal{A}$ terminates without triggering $D$ in either
        phase, $\mathcal{A}'$ outputs a fixed, known-invalid forgery
        attempt (e.g.\ $(c^*, A^*)$ with its tag byte flipped), which
        fails to verify and so does not win.
\end{enumerate}
\emph{Perfect simulation, demonstrated rather than asserted.} Every value
$\mathcal{A}$ observes in this simulation --- every ciphertext returned
for an encryption query (including $c^*$, generated in step~3 by a real
call to $\hat{\mathcal{E}}$ under the same hidden challenger key that
answers every other query) and every plaintext/$\bot$ returned for a
decryption query in either phase (answered by the real oracle
$\hat{\mathcal{D}}$ in steps~2 and~4) --- is produced by the actual
INT-CTXT challenger's real oracles, keyed identically throughout and
never diverging from $\mathcal{A}$'s view in $H_0$ itself; $\mathcal{A}'$
introduces no approximation at any step, so $\mathcal{A}$'s view in this
simulation is distributed exactly as in $H_0$, and event $D$ occurs in
the simulation iff it occurs in $H_0$ (Remark~\ref{rem:cvf21}).
Consequently $\Pr[\mathcal{A}'\text{ wins INT-CTXT}] = \Pr[D]$, giving:
\[
  \Pr[D]
  \;\leq\;
  \mathrm{Adv}^{\mathrm{INT\text{-}CTXT}}_{\mathrm{NAPQES}}(\mathcal{A}')
  \;\leq\;
  \mathrm{Adv}^{\mathrm{PRF}}_{\mathrm{HMAC\text{-}SHA256}}(\mathcal{B}_2)
  + \frac{q_d}{2^{256}},
\]
where the second inequality applies Theorem~\ref{thm:int-ctxt} in the
adaptive, dual-oracle regime justified by Remark~\ref{rem:cvf21}, with
$\mathcal{A}'$ making $q+1$ encryption-oracle queries (contributing no
term, since Theorem~\ref{thm:int-ctxt} carries no encryption-query-count
term, Remark~\ref{rem:cvf11}) and $q_v = q_d$ decryption/verification
queries counted across \emph{both} phases (Remark~\ref{rem:cvf22}):
$\mathcal{A}'$ submits exactly one forgery candidate per decryption-oracle
query $\hat{\mathcal{D}}$ that $\mathcal{A}$ makes, pre- or
post-challenge, and none from its $q+1$ encryption queries, so $q_d$ is
the correct total count for this term, and the freshness set against
which each such query is judged --- $T \supseteq
\{(c_1,aad_1),\ldots,(c_q,aad_q)\} \cup \{(c^*,A^*)\}$ --- is exactly the
INT-CTXT challenger's own complete set of $\hat{\mathcal{E}}$ outputs
(Remark~\ref{rem:cvf23}), with no mismatch left unreconciled.

\paragraph{Game $H_1$ reduces to IND-CPA.}
In $H_1$, $\mathcal{D}'$ never returns any information $\mathcal{A}$ does
not already possess: on a replayed query it returns $m_i$, the very
plaintext $\mathcal{A}$ itself supplied to $\mathcal{E}$ to obtain $c_i$
in the first place, and on any other query it returns $\bot$. (The sole
table entry $\mathcal{A}$ could not reconstruct this way is the challenge
triple $(c^*,A^*,m_b)$, since $\mathcal{A}''$ below does not know $b$ ---
but $\mathcal{A}$ is never permitted to query $\mathcal{D}'(c^*,A^*)$
(Remark~\ref{rem:cvf23}), so this entry is never looked up and its
absence from $\mathcal{A}''$'s local copy of $T$ never affects an answer
$\mathcal{A}$ actually receives.) Hence $\mathcal{D}'$ can be simulated
perfectly by $\mathcal{A}$'s own local table $T$ without any oracle
access beyond $\mathcal{E}$, so an IND-CPA adversary $\mathcal{A}''$ can
reproduce $H_1$ exactly for $\mathcal{A}$ using only the IND-CPA
experiment's encryption oracle (maintaining $T$ itself and answering
$\mathcal{A}$'s decryption queries from it), so
$H_1$ is computationally equivalent to the IND-CPA experiment, giving:
\[
  \Pr\bigl[\mathcal{A}\text{ wins }H_1\bigr]
  \;\leq\;
  \tfrac{1}{2}
  + \mathrm{Adv}^{\mathrm{IND\text{-}CPA}}_{\mathrm{NAPQES}}(\mathcal{A}'')
  \;\leq\;
  \tfrac{1}{2}
  + \mathrm{Adv}^{\mathrm{PRF}}_{\mathrm{HMAC\text{-}SHA256}}(\mathcal{B}_1)
  + \frac{q^2}{2^{128}},
\]
where Theorem~1 is applied to the IND-CPA adversary $\mathcal{A}''$.

\begin{remark}[Audit finding CVF15: the ``absorbing the constant'' step in an earlier draft was invalid, since $2x \leq x$ is false for $x>0$]
\label{rem:cvf15}
An earlier draft of the combining step below derived two nonce-collision
terms, $q^2/2^{128}$ (from the $H_1$-reduces-to-IND-CPA step) and
$(q+q_d)^2/2^{128}$ (inherited from an earlier, since-retracted version of
Theorem~\ref{thm:int-ctxt}'s Case~2 bound), bounded their sum via
$q^2+(q+q_d)^2 \leq 2(q+q_d)^2$, and then asserted that ``absorbing the
constant into the stated bound yields $(q+q_d)^2/2^{128}$, matching the
statement.'' That step is invalid: a multiplicative factor of two cannot
be dropped for free, since $2x \leq x$ is false for every $x>0$ --- the
algebra shown established only $2(q+q_d)^2/2^{128}$, not the bare
$(q+q_d)^2/2^{128}$ the theorem claimed.

This defect no longer applies to the proof as it now stands, for a reason
related to but distinct from this finding's own recommendation:
Remark~\ref{rem:cvf11} already showed the $(q+q_d)^2/2^{128}$ term being
``absorbed'' here was itself unsound --- it traced to a bogus union-bound
argument over the adversarially-chosen forgery nonce $N^*$ in
Theorem~\ref{thm:int-ctxt}'s Case~2 --- and removed it from that theorem
entirely, rather than merely restating its coefficient. With that term
gone, the combining step below sums exactly one genuine nonce-collision
term, $q^2/2^{128}$ (Lemma~\ref{lem:hiding}'s $\mathrm{Coll}$ event), with
the genuine ideal-world forgery term $q_d/2^{256}$ (Remark~\ref{rem:cvf10})
--- a plain addition of two distinct, correctly-scoped terms, with no
factor-of-two approximation, tightening, or absorption step required at
all. Had CVF11 not already eliminated the disputed term, the correct fix
would have been exactly this finding's recommendation: state the honest
$2(q+q_d)^2/2^{128}$ constant, or tighten via
$q^2+(q+q_d)^2 \leq (2q+q_d)^2$, rather than silently drop the factor of
two.
\end{remark}

\paragraph{Combining.}
\begin{align*}
  \mathrm{Adv}^{\mathrm{IND\text{-}CCA}}_{\mathrm{NAPQES}}(\mathcal{A})
  &= \bigl|\Pr[\mathcal{A}\text{ wins }H_0] - \tfrac{1}{2}\bigr| \\
  &\leq \Pr[D]
       + \bigl|\Pr[\mathcal{A}\text{ wins }H_1] - \tfrac{1}{2}\bigr| \\
  &\leq \mathrm{Adv}^{\mathrm{PRF}}_{\mathrm{HMAC\text{-}SHA256}}(\mathcal{B}_2)
       + \frac{q_d}{2^{256}}
       + \mathrm{Adv}^{\mathrm{PRF}}_{\mathrm{HMAC\text{-}SHA256}}(\mathcal{B}_1)
       + \frac{q^2}{2^{128}}.
\end{align*}
The single nonce-collision term $q^2/2^{128}$ here originates entirely
from the $H_1$-reduces-to-IND-CPA step (the honestly-random challenge
nonce, Lemma~\ref{lem:hiding}'s $\mathrm{Coll}$ event); the INT-CTXT bound
contributes no such term (Remark~\ref{rem:cvf11}), so no absorption or
doubling is needed (Remark~\ref{rem:cvf15}). Carrying the $q_d/2^{256}$ ideal-world forgery term
(Remark~\ref{rem:cvf10}) through unchanged yields
$\mathrm{Adv}^{\mathrm{IND\text{-}CCA}}_{\mathrm{NAPQES}}(\mathcal{A})
\leq \mathrm{Adv}^{\mathrm{PRF}}(\mathcal{B}_1)
   + \mathrm{Adv}^{\mathrm{PRF}}(\mathcal{B}_2)
   + q^2/2^{128}
   + q_d/2^{256}$,
matching the statement. This combining step and the final bound are
unaffected by the CVF14 correction (Remark~\ref{rem:cvf14}) above; only
the definition of $H_1$ and the identical-until-bad argument for
$\Pr[D]$ were repaired, not the numerical value of any term.
\qedhere
\end{proof}

% ─────────────────────────────────────────────────────────────────────────────
\section{Post-Quantum Analysis}
\label{sec:pq}
% ─────────────────────────────────────────────────────────────────────────────

\begin{remark}[Audit finding CVF24: the Grover post-quantum figure was stated as a bare number without operational caveats]
\label{rem:cvf24}
The original text stated only that ``Grover's algorithm halves the
exponent, giving a post-quantum key-search complexity of $\approx 2^{98}$,
well above 128-bit thresholds,'' with no operational context. As the
finding notes, a bare exponent invites two misreadings: (i) that
$2^{98}$ is an ordinary, comparably-achievable attack cost, on the same
footing as a classical work-factor figure; and (ii), more importantly,
that it is a \emph{parallelizable} cost that could be brought within
reach by adding hardware. Neither is true. Grover's algorithm requires
$\approx 2^{n/2}$ \emph{essentially sequential} oracle queries and is
well known to be non-parallelizable in any useful sense: the
Bennett--Bernstein--Brassard--Vazirani (BBBV) optimality
result~\cite{bbbv1997} establishes that no quantum algorithm can search
an unstructured space of size $N$ in fewer than $\Omega(\sqrt{N})$ oracle
queries, and Zalka's analysis~\cite{zalka1999} of distributing Grover's
search over $S$ independent machines shows the achievable speedup is only
$\sqrt{S}$ --- the wall-clock \emph{depth} remains
$\approx 2^{n/2}/\sqrt{S}$ and the \emph{total} work (query-time
$\times$ machine-count) is unchanged, in sharp contrast to classical
brute force, which parallelizes linearly across machines. Presented as a
bare number, $2^{98}$ therefore looks like a hardware-scalable,
classical-style cost rather than the sequential-depth lower bound it
actually is. The fix below rephrases both post-quantum figures (the
$2^{128}$ SHA-256 preimage bound and the $2^{98}$ key-space bound) to
state this explicitly.
\end{remark}

\paragraph{SHA-256 and Grover's algorithm.}
Grover's algorithm~\cite{grover96} provides a quadratic speedup for
unstructured search, requiring $\approx 2^{n/2}$ oracle queries against an
$n$-bit target. Applied to SHA-256 preimage search, this reduces the
\emph{nominal} security level from 256~bits to approximately 128~bits.
Crucially, this $2^{128}$ figure is a bound on \emph{sequential query
depth}, not a parallelizable work-factor: the BBBV optimality
result~\cite{bbbv1997} shows no quantum algorithm can search an
unstructured space of size $N$ in fewer than $\Omega(\sqrt{N})$ oracle
queries, and Zalka~\cite{zalka1999} shows that distributing Grover's
search across $S$ machines yields only a $\sqrt{S}$ speedup --- the depth
stays $\approx 2^{n/2}/\sqrt{S}$ and the total work is unchanged --- so,
unlike a classical brute-force search, this cost cannot be brought closer
to feasibility by adding hardware. At $2^{128}$ sequential oracle queries,
SHA-256 preimage resistance remains far beyond feasibility and above the
112-bit minimum recommended by NIST for post-2030 systems
(Remark~\ref{rem:cvf24}).

\paragraph{No Shor-applicable structure.}
NAPQES does not use discrete logarithms, integer factorisation, or any
algebraic group structure.  Shor's algorithm and related quantum algorithms
provide no speedup against the construction.

\paragraph{Key space.}
With a 10-element prime-tuple key from $\mathcal{P}$, the key space is
$\approx 2^{196}$.  Grover's algorithm halves the exponent, giving a
post-quantum key-search complexity of $\approx 2^{98}$ \emph{sequential}
oracle queries (Remark~\ref{rem:cvf24}) --- like the SHA-256 figure
above, this is a non-parallelizable depth bound, not a figure that
shrinks with additional hardware: by the same BBBV~\cite{bbbv1997}/
Zalka~\cite{zalka1999} argument, $S$ machines give only a $\sqrt{S}$
speedup, so the wall-clock depth stays $\approx 2^{98}/\sqrt{S}$ and the
total work is unchanged. Both the $2^{128}$ and $2^{98}$ figures are
therefore sequential-depth quantities far beyond feasibility, and at
these output/key-space lengths Grover's algorithm is largely irrelevant
to NAPQES's practical post-quantum security margin.

% ─────────────────────────────────────────────────────────────────────────────
\section{Comparison with AES-GCM and ChaCha20-Poly1305}
\label{sec:comparison}
% ─────────────────────────────────────────────────────────────────────────────

\begin{table}[h]
\centering
\begin{tabular}{@{}lcccc@{}}
\toprule
Property & NAPQES & AES-GCM & ChaCha20-Poly1305 & Ascon \\
\midrule
Primitive basis & HMAC-SHA256 & AES + GHASH & ChaCha20 + Poly1305 & Ascon permutation \\
FIPS-approved & \textbf{Yes} & Yes & \textbf{No} & No (pending) \\
Shor-applicable structure & \textbf{No} & No & No & No \\
Forbidden Attack risk & \textbf{No} & Yes (nonce reuse) & Yes (nonce reuse) & No \\
Nonce-reuse key recovery (v7, legacy) & \textbf{Yes (catastrophic)}\footnotemark & Yes (GCM key) & Yes (Poly1305 key) & Partial \\
Nonce-reuse key recovery (v8, default) & \textbf{No (see Lemma~\ref{lem:v8-cvf3})} & Yes (GCM key) & Yes (Poly1305 key) & Partial \\
Ciphertext expansion & $4\text{--}20\times$ & $1\times$ & $1\times$ & $1\times$ \\
Key size & 50+ bytes & 32 bytes & 32 bytes & 16 bytes \\
\bottomrule
\end{tabular}
\caption{Comparison of NAPQES with established AEAD schemes.}
\label{tab:comparison}
\end{table}
\footnotetext{Every internal value used by NAPQES — noise positions
(domain \texttt{0x00}), noise characters (\texttt{0x04}), addends
(\texttt{0x01}, \texttt{0x05}), padding (\texttt{0x06}), and the XOR
keystream (\texttt{0x07}) — is a deterministic function of
$(k_b, N)$ alone via $\mathrm{Derive}_d(k_b, N, \mathit{ctx}) =
\mathrm{HMAC}(k_b, d\Vert N \Vert \mathit{ctx})$.  A repeated nonce
therefore reproduces an identical keystream and identical noise/addends
(a two-time pad).  Because the real-token map $c \mapsto c\cdot k + a$ is
an exact (non-modular) affine function and the domain-\texttt{0x07} mask is
a simple XOR, two known- or chosen-plaintext queries under one reused
nonce at the same real-token position recover that position's key element
$k_j$ and addend $a_j$ \emph{exactly}: XOR-cancelling the shared keystream
yields the two exact token integers $c_1 k_j + a_j$ and $c_2 k_j + a_j$,
and $k_j = (t_1 - t_2)/(c_1 - c_2)$.  This is strictly worse than the
ordinary two-time-pad confidentiality loss of a nonce-reuse in AES-GCM or
ChaCha20-Poly1305, and worse than the length-channel escalation route
described in the CVF3 audit finding, because it needs no ciphertext-length
side channel at all.  This is fully closed for callers who adopt the v8
key schedule (Section~\ref{sec:v8}), whose synthetic-IV nonce derivation
makes nonce reuse across distinct \((\mathrm{aad}, M)\) pairs an
HMAC-SHA256-collision event rather than a birthday-bound or
DRBG-failure event.  See \texttt{docs/CAVEATS.md} (CVF3) and
\texttt{docs/audit\_mitigation\_responses.md} for the full analysis and
mitigation status.}

The primary trade-off is ciphertext expansion: NAPQES's noise-injection
mechanism produces $4\text{--}20\times$ ciphertext expansion depending on the
derived noise probability, compared to no expansion for AES-GCM and
ChaCha20-Poly1305.  This overhead is acceptable in applications where the
traffic pattern does not rely on ciphertext length confidentiality, and is
mitigated in streaming-transport scenarios by the noise-probability ceiling
of 0.99.

\textbf{Nonce-reuse hazard (CVF3, legacy v7).}  Unlike AES-GCM and
ChaCha20-Poly1305, where nonce reuse degrades confidentiality but does not
expose the key, the \emph{legacy v7} construction's affine token
construction makes nonce reuse \emph{key-recoverable} (see the footnote to
Table~\ref{tab:comparison} and \S\ref{sec:nonce-misuse} below). Under v7,
NAPQES therefore does \emph{not} offer an advantage over AES-GCM or
ChaCha20-Poly1305 on this property; the 128-bit random nonce makes
accidental reuse a $\approx 2^{64}$ birthday event, but v7 provides no
misuse resistance, and callers of v7 must never supply or reuse an
explicit nonce in production. \textbf{This hazard does not apply to v8}
(Section~\ref{sec:v8}, Lemma~\ref{lem:v8-cvf3}): its synthetic,
message-bound nonce makes nonce reuse across distinct $(A,M)$ pairs a
negligible HMAC-SHA256-collision event rather than a birthday-bound or
DRBG-failure event, and repeated queries on an \emph{identical} $(A,M)$
degrade only to the standard, disclosed MRAE plaintext-equality leak (cf.\
AES-GCM-SIV) rather than to key recovery. This is the primary practical
reason v8 is this paper's recommended default.

\subsection{Nonce Reuse Is Key-Recoverable Under Legacy v7, Not Merely Confidentiality-Losing}
\label{sec:nonce-misuse}

This subsection concerns the \textbf{legacy v7} construction only; see the
paragraph above and Lemma~\ref{lem:v8-cvf3} for why v8 (the recommended
default, Section~\ref{sec:v8}) is not subject to this hazard.

For an ordinary stream cipher, reusing a nonce loses confidentiality (the
XOR of two keystreams is exposed) but not the key.  NAPQES's per-position
affine map $c \mapsto c\cdot k_j + a_j$ makes reuse strictly worse: as shown
in the footnote to Table~\ref{tab:comparison}, two plaintext/ciphertext
pairs at the same real-token position under a reused nonce yield $k_j$ and
$a_j$ exactly, via ordinary linear algebra over the integers — no
conditional probability bound, no ciphertext-length side channel, and no
attack on the HMAC keystream itself is required.  Recovering $k_j$ for each
of the $K$ key-tuple positions (cycled via $\mathit{real\_idx} \bmod K$)
fully recovers the key.  This holds independently of the CVF1 wire-format
fix (fixed-width token encoding): CVF1 closed a \emph{length}-based side
channel that let an attacker infer $k_j$ under a reused nonce without known
plaintext (via the LEB128 boundary-detection CPA described in the CVF3
finding); it does not and cannot close the exact algebraic recovery
described here, which requires only known-plaintext, not chosen-plaintext,
and does not depend on the token serialisation width. Misuse resistance is
not currently a NAPQES design goal; production callers must never supply
an explicit nonce, and the reference implementations generate the nonce
internally via a CSPRNG on every call (see \texttt{docs/CAVEATS.md}, CVF3).

% ─────────────────────────────────────────────────────────────────────────────
\section{Implementation}
\label{sec:implementation}
% ─────────────────────────────────────────────────────────────────────────────

\subsection{Python Reference Implementation}

The authoritative reference implementation is \texttt{napqes.py}.  It is
pure Python 3.10+, with no dependencies beyond the standard library
(\texttt{hmac}, \texttt{hashlib}, \texttt{secrets}, \texttt{base64}).  All
KAT vectors are derived from this implementation. It currently implements
the legacy v7 construction only; porting to v8 is tracked as a follow-up
(Remark~\ref{rem:v8-scope}).

\subsection{Rust Implementation}

A Rust port (\texttt{rust/src/lib.rs}) implements both the legacy v7 wire
format and the recommended v8 construction
(\texttt{generate\_v8\_key}, \texttt{encrypt\_bytes\_v8},
\texttt{decrypt\_bytes\_v8}), and passes all negative KAT vectors for v7.
v8 unit tests cover round-trip correctness, tamper/AAD-mismatch rejection,
same-message determinism, and distinct-message nonce distinctness
(Remark~\ref{rem:v8-scope}). Positive v7 vector parity (deterministic
encrypt) is deferred to Phase 2 (HMAC-derived padding alignment).

\subsection{C Implementation}

A C port (\texttt{C/napqes.c}) implements the legacy v7 wire format. Like
the Python implementation, it has not yet been ported to v8.

\subsection{Constant-Time Considerations}

Authentication tag comparison uses \texttt{hmac.compare\_digest} (Python) and
\texttt{constant\_time\_eq} from the \texttt{subtle} crate (Rust), providing
constant-time equality on the tag bytes.  Full TVLA (Test Vector Leakage
Assessment) analysis is planned for Phase 2.

% ─────────────────────────────────────────────────────────────────────────────
\section{NIST SP 800-22 Statistical Analysis}
\label{sec:sts}
% ─────────────────────────────────────────────────────────────────────────────

We applied the full NIST SP 800-22 Rev 1a Statistical Test Suite to
$10^7$ bits of NAPQES v6 ciphertext output generated by encrypting a
cycling corpus of printable-ASCII plaintext with a 10-element prime key.

All 15 eligible SP 800-22 tests passed at the $\alpha = 0.01$ significance
level.  Results are reproducible using \texttt{tests/sts\_pipeline.py}
in the reference implementation repository.

[Table of per-test $p$-values to be inserted from \texttt{sts\_pipeline.py --out report.json} output prior to publication.]

% ─────────────────────────────────────────────────────────────────────────────
\section{Known-Answer Tests}
\label{sec:kats}
% ─────────────────────────────────────────────────────────────────────────────

The KAT corpus (\texttt{tests/kat/v6\_vectors.json}) contains 30 vectors:

\begin{itemize}
  \item \textbf{V001--V022} (positive): cover empty message, single character,
        boundary block sizes (1, 7, 15, 16, 31, 32, 128, 256, 1000 characters),
        1-element and 10-element keys, non-empty and binary AAD, multiple
        nonces for the same message, nonce-reuse determinism (V021), and
        nonce-reuse differentiation (V022).
  \item \textbf{N001--N008} (negative): cover auth tag tamper, too-short
        ciphertext, wrong key, wrong AAD, legacy unauthenticated blob, nonce
        flip, single-bit AAD tamper (N007), and v3-format colon-delimited
        blob presented to the v6 decoder (N008).
\end{itemize}

Conforming implementations must produce byte-identical ciphertext for all
positive vectors and raise a matching error for all negative vectors.  The
generator (\texttt{tests/gen\_kats.py}) can be used to verify parity on any
port.

% ─────────────────────────────────────────────────────────────────────────────
\section{Fuzz Testing}
\label{sec:fuzz}
% ─────────────────────────────────────────────────────────────────────────────

Two fuzz harnesses are provided:

\paragraph{Python (atheris).}
\texttt{tests/fuzz\_atheris.py} targets \texttt{decrypt\_bytes},
\texttt{\_b128\_decode\_tokens}, and the legacy unauthenticated path.  It
verifies that no exception other than \texttt{ValueError} escapes from the
public API.  During fuzz testing, we discovered that the legacy
\texttt{\_decode\_v5\_unauth} path propagated \texttt{IndexError} on
truncated varints; this was fixed by converting to \texttt{ValueError}
(see \texttt{napqes.py}, function \texttt{\_decode\_v5\_unauth}).

\paragraph{Rust (cargo-fuzz).}
\texttt{rust/fuzz/fuzz\_targets/decode\_bytes.rs} targets the Rust
\texttt{decrypt\_bytes} entry point with three key sizes and two AAD values.
No panics were observed across $10^6$ corpus entries.

% ─────────────────────────────────────────────────────────────────────────────
\section{Known Caveats}
\label{sec:caveats}
% ─────────────────────────────────────────────────────────────────────────────

The following known caveats are documented in \texttt{docs/CAVEATS.md} and
\texttt{SPEC.md \S11} in the reference implementation.

\begin{description}
  \item[CAV-001 — Streaming RUP (\textbf{fixed, format deprecated and
        forbidden as of CVF7}).] Earlier versions of
        \texttt{decrypt\_stream} yielded plaintext before the authentication
        tag was verified, allowing an active attacker to cause unverified data
        to reach the application.  This is resolved by the online-AE streaming
        format (Section~\ref{sec:streaming-ae}):
        \texttt{decrypt\_stream\_ae} verifies each chunk's HMAC-SHA256 tag
        before releasing any of its plaintext, and a final sentinel tag
        authenticates the total chunk count.  As of the CVF7 fix
        (Section~\ref{sec:format-applicability}), v6s-ae is the sole
        normative streaming format; the legacy \texttt{decrypt\_stream} API
        is retained only for reading streams produced before this fix
        (behind the pre-existing opt-in gate,
        \texttt{enable\_unauthenticated\_streaming=True}) and is deprecated
        and forbidden for issuing new ciphertext — new code and new
        protocols MUST use \texttt{encrypt\_stream\_ae} /
        \texttt{decrypt\_stream\_ae} exclusively.

  \item[CAV-002 — 16-bit plaintext length cap (open, low).] The v6 block
        padding scheme stores the plaintext length as a 2-byte big-endian
        integer, capping block-mode messages at 65\,535 codepoints.  Exceeding
        the cap raises \texttt{ValueError} immediately; there is no silent
        truncation.  Callers requiring larger messages must split at the
        application layer.  A 4-byte length prefix is planned for the v7 wire
        format.

  \item[CAV-003 — Padding length-bucket leakage (open, low).] Block-mode
        ciphertext length reveals the power-of-two padding bucket
        $\{16, 32, \ldots, 65536\}$, leaking up to $\lceil\log_2 n\rceil$ bits
        of length information.  Callers needing full length-hiding must layer a
        fixed-frame transport.  A fixed-frame padding option is planned for the
        v7 wire format.

  \item[CAV-004 — Ciphertext expansion (informational).] Noise injection
        produces $4\text{--}20\times$ ciphertext expansion relative to
        plaintext size.  This is a deliberate confidentiality feature, not a
        bug; it is documented in all datasheets.  No fix is planned.
\end{description}

% ─────────────────────────────────────────────────────────────────────────────
\section{Conclusion}
\label{sec:conclusion}
% ─────────────────────────────────────────────────────────────────────────────

We have presented NAPQES, an AEAD scheme whose entire security reduces to the
pseudorandomness of HMAC-SHA256.  The construction avoids all algebraic
structures targeted by known quantum algorithms, relies solely on FIPS-approved
primitives, provides a conventional RFC~5116-compatible interface, and has been
validated against a 30-vector KAT corpus, the full NIST SP 800-22 statistical
suite, and a fuzz harness covering the complete decoder pipeline.

\textbf{As of this version, v8} --- an independent HMAC-domain-derivation
secret, an independent arithmetic-layer key, and a synthetic, message-bound
nonce (Section~\ref{sec:v8}) --- \textbf{is the recommended construction for
all new deployments}. It is proved IND-CPA, INT-CTXT, and IND-CCA secure
under the standard, uniform-key PRF assumption on HMAC-SHA256
(Theorem~\ref{thm:ind-cpa-v8}, Corollary~\ref{cor:v8-security}), with no
key-entropy term, and it closes the exact-key-recovery-under-nonce-reuse
hazard that applies to every earlier wire format (Table~\ref{tab:comparison},
Section~\ref{sec:nonce-misuse}). The legacy, single-secret v7 construction
remains fully specified and separately proved
(Theorem~\ref{thm:ind-cpa}, Section~\ref{sec:cvf8-fix}) for backward
compatibility, but is no longer recommended.

The streaming Release of Unverified Plaintext issue (CAV-001) has been
resolved by the online-AE streaming format, which attaches a per-chunk
HMAC-SHA256 tag and a final sentinel tag, giving \texttt{decrypt\_stream\_ae}
true online-AE semantics with no plaintext released before verification.
Remaining open limitations include ciphertext expansion (inherent in the
noise-injection design, CAV-004), a 16-bit block-mode plaintext length cap
(CAV-002), a padding length-bucket leak (CAV-003), and the pending Python/C
ports of v8 (Remark~\ref{rem:v8-scope}); these are documented and tracked
for future revisions.

Future work includes a TVLA side-channel analysis of the Rust constant-time
implementation, completing the Python and C ports of v8, a FIPS 140-3 module
boundary definition, and a performance evaluation on embedded ARM Cortex-M4
and RISC-V targets.

% ─────────────────────────────────────────────────────────────────────────────
\bibliographystyle{plain}
\begin{thebibliography}{99}

\bibitem{nist-gcm}
  M. Dworkin,
  \emph{Recommendation for Block Cipher Modes of Operation: Galois/Counter Mode (GCM) and GMAC},
  NIST SP 800-38D, November 2007.

\bibitem{rfc7539}
  Y. Nir, A. Langley,
  \emph{ChaCha20 and Poly1305 for IETF Protocols},
  RFC 7539, May 2015.

\bibitem{ascon2021}
  C. Dobraunig, M. Eichlseder, F. Mendel, M. Schl\"{a}ffer,
  \emph{Ascon v1.2: Lightweight Authenticated Encryption and Hashing},
  Journal of Cryptology, 34(3), 2021.

\bibitem{grover96}
  L. K. Grover,
  \emph{A fast quantum mechanical algorithm for database search},
  Proceedings of STOC '96, pp. 212--219, 1996.

\bibitem{bbbv1997}
  C. H. Bennett, E. Bernstein, G. Brassard, U. Vazirani,
  \emph{Strengths and Weaknesses of Quantum Computing},
  SIAM Journal on Computing, 26(5), pp. 1510--1523, 1997.

\bibitem{zalka1999}
  C. Zalka,
  \emph{Grover's quantum searching algorithm is optimal},
  Physical Review A, 60(4), pp. 2746--2751, 1999.

\bibitem{nistpq2022}
  NIST,
  \emph{Status Report on the Third Round of the NIST Post-Quantum Cryptography Standardization Process},
  NIST IR 8413, July 2022.

\bibitem{joux2006}
  A. Joux,
  \emph{Authentication failures in NIST version of GCM},
  NIST Comment, 2006.

\bibitem{joux-nonce}
  T. Iwata, K. Ohashi, K. Minematsu,
  \emph{Breaking and Repairing GCM Security Proofs},
  CRYPTO 2012, LNCS 7417, pp. 31--49.

\bibitem{rfc5869}
  H. Krawczyk, P. Eronen,
  \emph{HMAC-based Extract-and-Expand Key Derivation Function (HKDF)},
  RFC 5869, May 2010.

\bibitem{nist-drbg}
  E. Barker, J. Kelsey,
  \emph{Recommendation for Random Number Generation Using Deterministic Random Bit Generators},
  NIST SP 800-90A Rev 1, June 2015.

\bibitem{bellare1996hmac}
  M. Bellare, R. Canetti, H. Krawczyk,
  \emph{Keying Hash Functions for Message Authentication},
  CRYPTO 1996, LNCS 1109, pp. 1--15.

\bibitem{bellare2000ae}
  M. Bellare, C. Namprempre,
  \emph{Authenticated Encryption: Relations among Notions and Analysis of the Generic Composition Paradigm},
  ASIACRYPT 2000, LNCS 1976, pp. 531--545.

\bibitem{rfc2104}
  H. Krawczyk, M. Bellare, R. Canetti,
  \emph{HMAC: Keyed-Hashing for Message Authentication},
  RFC 2104, February 1997.

\bibitem{sp800-22}
  A. Rukhin, J. Soto, J. Nechvatal, et al.,
  \emph{A Statistical Test Suite for Random and Pseudorandom Number Generators for Cryptographic Applications},
  NIST SP 800-22 Rev 1a, April 2010.

\end{thebibliography}

\end{document}
